\documentclass[11pt,a4paper]{article}
\usepackage[T1]{fontenc}
\usepackage[utf8]{inputenc}
\usepackage{lmodern}
\usepackage{microtype}
\usepackage[margin=25mm,headheight=15pt]{geometry}
\usepackage{amsmath,amssymb,amsthm,mathtools}
\usepackage{booktabs,longtable,tabularx,array}
\usepackage{xcolor}
\usepackage{enumitem}
\usepackage{listings}
\usepackage{fancyhdr}
\usepackage{titlesec}
\usepackage{hyperref}
\usepackage{bookmark}
\usepackage{xurl}
\definecolor{clover}{HTML}{315A45}
\definecolor{pale}{HTML}{F2F5F1}
\definecolor{ink}{HTML}{202822}
\definecolor{muted}{HTML}{59655D}
\hypersetup{colorlinks=true,linkcolor=clover,citecolor=clover,urlcolor=clover,
 pdftitle={Clover: Scoped Mathematical Reasoning with Obligation-Directed Types},
 pdfauthor={GPT-6 Astra Pro},pdfsubject={Round-four proof-language proposal for Vladimir Reshetnikov}}
\titleformat{\section}{\Large\bfseries\color{clover}}{\thesection}{.65em}{}
\titleformat{\subsection}{\large\bfseries}{\thesubsection}{.65em}{}
\titleformat{\subsubsection}{\normalsize\bfseries}{\thesubsubsection}{.65em}{}
\setlength{\parindent}{0pt}
\setlength{\parskip}{.47em}
\setlength{\emergencystretch}{2em}
\setcounter{tocdepth}{2}
\setlist{itemsep=.2em,topsep=.4em,leftmargin=1.6em}
\pagestyle{fancy}
\fancyhf{}
\fancyhead[L]{\small\color{muted}Clover \;\textbar\; Round-four proposal}
\fancyhead[R]{\small\color{muted}\nouppercase{\leftmark}}
\fancyfoot[C]{\small\thepage}
\renewcommand{\headrulewidth}{.3pt}
\renewcommand{\sectionmark}[1]{\markboth{#1}{}}
\lstdefinestyle{source}{basicstyle=\ttfamily\small,columns=fullflexible,
 keepspaces=true,showstringspaces=false,breaklines=true,breakatwhitespace=false,
 frame=single,framerule=.25pt,rulecolor=\color{clover},backgroundcolor=\color{pale},
 xleftmargin=.3em,xrightmargin=.3em,aboveskip=.7em,belowskip=.7em,
 captionpos=b,tabsize=2}
\lstset{style=source}
\newcommand{\code}[1]{\texttt{\detokenize{#1}}}
\newcommand{\N}{\mathbb N}
\newcommand{\Z}{\mathbb Z}
\newcommand{\Q}{\mathbb Q}
\newcommand{\R}{\mathbb R}
\newcommand{\Prop}{\mathsf{Prop}}
\newcommand{\Type}{\mathsf{Type}}
\newcommand{\Beh}{\mathsf{Beh}}
\newcommand{\CRef}{\mathsf{Ref}}
\newcommand{\id}{\mathrm{id}}
\newcommand{\den}[1]{\llbracket #1\rrbracket}
\newcommand{\AEeq}{\mathrel{=_{\mathrm{a.e.}}}}
\newcolumntype{Y}{>{\raggedright\arraybackslash}X}
\newcolumntype{P}[1]{>{\raggedright\arraybackslash}p{#1}}
\newtheorem{theorem}{Theorem}[section]
\newtheorem{proposition}[theorem]{Proposition}
\newtheorem{lemma}[theorem]{Lemma}
\newtheorem{corollary}[theorem]{Corollary}
\theoremstyle{definition}
\newtheorem{definition}[theorem]{Definition}
\theoremstyle{remark}
\newtheorem{remark}[theorem]{Remark}
\newenvironment{designnote}{\begin{quote}\small\color{ink}}{\end{quote}}

\begin{document}
\begin{titlepage}
\vspace*{16mm}
{\color{clover}\rule{\textwidth}{1.1pt}}
\vspace{12mm}
{\fontsize{42}{46}\selectfont\bfseries\color{clover}Clover\par}
\vspace{6mm}
{\LARGE\bfseries Scoped Mathematical Reasoning\par}
\vspace{2mm}
{\LARGE\bfseries with Obligation-Directed Types\par}
\vspace{9mm}
{\large A proof-language proposal for the fourth\par
iteration of the Brainstorm campaign\par}
\vspace{14mm}
{\large Prepared for Vladimir Reshetnikov\par}
{\normalsize GPT-6 Astra Pro\par 6 September 2026\par}
\vfill
\begin{minipage}{.94\textwidth}\setlength{\parskip}{.55em}
\textbf{Central proposal.} Extend the author-facing typing discipline with
scoped properties, quantified evidence, and dependent proof obligations.
Make common mathematical moves into typed proof transformations.
Resolve their routine premises by a bounded, proof-producing engine; keep
Lean's kernel conversion unchanged in the first implementation.

\textbf{Concrete contribution.} A specified grounding algorithm; a restricted
executable source-to-proof-plan-to-Lean pipeline; a realizable affine test-basis
theorem for behavioral synthesis; and explicit protocols for unfinished
claims, meaning preservation, and scope-safe replay.

\textbf{Evidence status.} The Python companion was executed. The Lean reference
and generated Lean examples were not compiled in this environment. The full
Clover surface, editor integration, and proposed performance study remain designs.
\end{minipage}\par
\vspace{10mm}
{\color{clover}\rule{\textwidth}{.5pt}}
\end{titlepage}

\begin{abstract}
Mathematicians usually state the transformation that matters---differentiate an
identity on an open set, transport an exact quotient, restrict a construction
to an invariant domain---and leave its routine premises implicit. Formal proof
languages must nevertheless account for those premises, preserve the identity
of the objects involved, and distinguish operations whose familiar notation
conceals different semantics. Clover proposes an author-facing extension of
Lean's typing discipline that treats this omitted work as a typed, scoped
obligation problem rather than as an invitation to unrestricted proof search.

The design starts from both revised round-three syntheses, not from a fresh
catalogue of features. It retains their conservative foundations, then makes
four interfaces more concrete: quantified evidence and mathematical proof
transformations; unfinished elaboration represented by dependent obligation
telescopes; finite, goal-triggered rule generation with exact failure states;
and behavioral abstractions equipped with realizability and coverage data.
A test-basis theorem extends the usual affine-length criterion to correlated
observations and restricted input domains. Worked examples preserve the precise
hypotheses of the autonomous-derivative and Frey-curve proofs, explain naturality
in a finite tensor construction, and expose a quantifier-order trap in
almost-everywhere reasoning.

A small companion actually parses a restricted language, generates and checks
proof plans, and exports ordinary Lean source. Its 35 regression tests pass;
a separate exact-arithmetic companion performs 1,576 checks. These results do
not establish a verified Lean frontend or a productivity improvement. They do
expose a concrete inference-cost issue and demonstrate the reduction obtained
by a narrower rule profile. The article concludes with a serious comparison
of elaborator-level refinements and possible kernel extensions, a staged
implementation plan, and an evaluation designed to distinguish language benefits
from improvements already obtainable through libraries and automation.
\end{abstract}
\clearpage
\tableofcontents
\clearpage

\section{The proposal and its boundary}\label{sec:proposal}

Clover is a language for \emph{specified mathematical reasoning steps}. Its
smallest useful unit is not a sentence accepted because an AI found some proof
nearby. It is a step with a fixed meaning, a typed conclusion, explicit local
scope, and a checked account of the premises that justify it.

The author should usually write the mathematical distinction that selects the
argument, not every fact that follows routinely from that distinction. In
particular, the source should retain that an identity holds on an open region,
that a quotient is exact in the integers before it is embedded into a field,
or that an equivalence is natural in its target. It should not have to repeat
neighborhood construction, denominator casts, and componentwise congruence at
every use. The implementation should discover or assemble those details,
without weakening the statement or silently choosing a different mathematical
structure.

The extension has two levels. At the author-facing level, an expression has
an underlying Lean type, an interface of proved properties, and possibly an
unfinished obligation telescope. At the kernel level, accepted expressions
remain ordinary Lean terms with ordinary dependent types. Thus Clover extends
what the elaborator \emph{uses when checking a program or proof}; it does not
initially extend the kernel's notion of definitional equality.

This distinction is substantive. A conventional elaborator asks whether a
term inhabits a type and resolves implicit arguments. Clover additionally asks
which mathematical requirements a selected operation imposes, which scoped
facts discharge them, which theorem-backed adapters are needed, and what
remains unfinished. Its user-visible types include behavioral contracts even
when the corresponding Lean function retains a familiar unrefined carrier
type.

\subsection{Four design commitments}

\textbf{Preserve meaning before searching for evidence.} Objects, quantifier
order, structures, domains, measures, and requested conclusion strength must
be determined independently of whether a convenient proof is found. An
authorized construction is different: its witness may be chosen, but its
specification is fixed.

\textbf{Make the mathematical move explicit.} A source step such as
\code{differentiate on U} selects a proof transformation with a precise
interface. A source step such as \code{search} delegates the route. These are
not two spellings of the same operation.

\textbf{Treat pending conditions as pending.} An elaborator may compute a
conditional proof while some guards remain open. That conditional proof must
not enter the evidence index as an unconditional fact, and an unfinished
intermediate assertion must not disappear merely because the final theorem
has another proof.

\textbf{Bound routine inference; expose escape routes.} The default engine
works in a finite source-derived term arena and a selected theorem profile.
When it cannot close an obligation, the author may supply a lemma, choose a
larger profile, introduce a missing construction, or enter ordinary Lean.
Failure of this engine is not mathematical falsity.

\subsection{What is new here, and what is inherited}

The round-three campaign already established most of the underlying refinement
calculus. Calling the same calculus Clover would not by itself advance the
project. The proposed increment is therefore operational and interface-level:

\begin{center}\small
\begin{tabularx}{\textwidth}{@{}P{.25\textwidth}Y Y@{}}
\toprule
Area & Inherited foundation & Clover's additional specification\\\midrule
Properties & Evidence about a fixed object & Quantified evidence schemes and their controlled specialization.\\
Proof work & A frontier of missing premises & Dependent obligation telescopes with explicit pending/closed transitions.\\
Inference & Finite ground closure & A source-derived arena, conclusion-triggered instantiation, and bounded premise-only parameters.\\
Behavior & Sound abstraction and realizable refutation & Affine test bases for the actual observation image, including correlations and guards.\\
Implementation & Small semantic cores & An executed restricted parser, proof-plan generator, independent checker, and Lean exporter.\\
\bottomrule
\end{tabularx}
\end{center}

These contributions have different evidence levels. The calculus and general
algorithms are specified and justified mathematically in the article. The
endomorphism slice is executed Python. Its ordinary-Lean rule reference is
written but not freshly kernel-checked here. Nothing in this package should
be described as a completed general-purpose Clover compiler.

\section{Reading the third round as an implementation brief}\label{sec:round3}

Both revised syntheses were read at Brainstorm commit
\code{bf3333905995}:
\emph{From Properties to Proof Obligations} and \emph{Nine Refinements}
\cite{syn-a,syn-b}. They agree much more than their different tables suggest.
One organizes ten architectural commitments; the other scores a finer-grained
set of commitments and collects a larger mutation catalogue. These are editorial
views of related proposals, not independent measurements of language quality.

Their strongest recommendation is to implement a small interface rather than
produce another overlapping design. Clover accepts that recommendation in a
limited but concrete way: the companion supplies an executable operational
slice, while the article specifies the parts that must be added to make it a
real Lean extension. It does not present the small model as having completed
the exact-quotient or analytic integration tasks requested by the syntheses.

\subsection{The inherited distinctions must survive compression}

The reports separate \emph{object}, \emph{refinement}, \emph{observation},
\emph{construction}, and \emph{operational receipt}. An object is a fixed term
with its selected structures. A refinement is more evidence about that term.
An observation is related information, such as a prefix, enclosure, or quotient
image. A construction produces a potentially different witness. A receipt
records what an implementation did; it is not automatically a proof of its
mathematical claim. Clover preserves all five distinctions.

Several earlier mistakes show why they matter. A finite observation can support
a universal proof through a sound theorem without making the method complete.
A failing abstract length is not a source-level counterexample until it is
realized by an admissible input. Normalizing a counterexample creates a new
witness rather than another property of the old witness. A finite index set
does not make its algebra factors finite as modules. A final contradiction
can prove an arithmetic helper while bypassing the arithmetic argument that a
benchmark intended to exercise. The revised syntheses explicitly correct these
confusions \cite{syn-a,syn-b}.

\subsection{What Clover takes from the individual proposals}

The following attribution follows the syntheses' detailed discussion; it is
not a claim of a new audit of every original proposal and companion.

\begin{center}\small
\begin{longtable}{@{}P{.18\textwidth}P{.73\textwidth}@{}}
\toprule
Proposal & Design lesson retained in Clover\\\midrule\endfirsthead
\toprule Proposal & Design lesson retained in Clover\\\midrule\endhead
\bottomrule\endfoot
Basalt & Properties of diagrams and coupled objects must remain relational; mathematical structures and null-boundary conditions are part of the proposition.\\
Fiber & Exact quotients should be tested under satisfiable, operation-specific assumptions; synthesis holes need contracts on the actual intermediate values.\\
Gneiss & Infinite objects may admit lawful finite observations only because an operation preserves the relevant locality.\\
Karst & Evidence is scope-sensitive; proposition-valued finiteness is not a chosen enumeration; \code{List Empty} tests negative realizability.\\
Moraine & A restricted syntax can support a universal behavior theorem through an interpreter and an affine model; naturality is not a fieldwise annotation.\\
Obsidian & Normalization changes witnesses; existing Lean facilities are serious controls; an inconsistent ambient package is a poor arithmetic benchmark.\\
Schist & Backward requirements express the way a selected mathematical method asks for premises; small complete checking chains are valuable.\\
Tephra & Transfer, coverage, and realization are different theorems; arbitrary polymorphic Lean definitions do not receive parametricity for free.\\
Trellis & Closure is predictable only after the finite inference fragment is fixed; bound quality and sharpness are different contracts.\\
\end{longtable}
\end{center}

\subsection{Three gaps are more urgent than new notation}

First, a finite closure theorem leaves rule generation unspecified. A practical
system needs to say which theorem parameters it instantiates, where their
values come from, and when enumeration stops. Section~\ref{sec:grounding}
answers that question for a conservative fragment and implements it for a
smaller one.

Second, a fixed-object evidence row is incomplete as an authoring model unless
it retains \emph{quantified scope}. The autonomous-derivative proof works with
an induction hypothesis valid throughout a region, not merely at its current
point. Section~\ref{sec:quantified} treats this as an interface requirement.

Third, source fidelity needs an actual protocol. A provider must not solve a
statement parameter while claiming to solve a proof hole. A stale local name
must not identify an object after a declaration has been rebuilt. Sections
\ref{sec:meaning} and \ref{sec:transport} specify these boundaries separately
from theorem truth.

\section{What the inspected Lean code actually teaches}\label{sec:sources}

The source studies below are development examples, not held-out benchmarks.
They were inspected through immutable GitHub references; the repositories were
not built or executed in preparing this article. The machine-readable source
manifest records the exact paths and scope of review.

\subsection{An autonomous differential equation: locality and quantifiers}

In ProveIt's \code{AutonomousIteratedDeriv.lean}, the theorem
\code{iteratedDeriv_eq_comp} proves a general mechanism behind
higher derivatives \cite{autonomous}. Let $U\subseteq\R$ be open, let
$W,\varphi:\R\to\R$, and let $G_n,G'_n:\R\to\R$. The hypotheses assert,
only at the points actually needed, that
\begin{align}
 W'(x)&=\varphi(W(x)) &&(x\in U),\label{eq:ode}\\
 G_0(W(x))&=W(x) &&(x\in U),\label{eq:gzero}\\
 \frac{d}{dw}G_n(w)\bigg|_{w=W(x)}&=G'_n(W(x)) &&(x\in U),\label{eq:gderiv}\\
 G_{n+1}(W(x))&=G'_n(W(x))\varphi(W(x)) &&(x\in U).\label{eq:grec}
\end{align}
The derivative hypotheses here mean \emph{existence with the specified value},
as expressed by \code{HasDerivAt} in the source. They are stronger than a bare
equality involving Lean's totalized derivative function.

The conclusion is
\[
 \forall n\in\N\;\forall x\in U,\qquad
 D^nW(x)=G_n(W(x)).
\]
Its successor step first turns the regional induction hypothesis into equality
in a neighborhood of $x$, then differentiates there and applies the chain
rule. In the source, the key bridge is the application of
\code{Filter.eventuallyEq_of_mem} to openness, membership, and the
\emph{unspecialized} induction hypothesis.

Two details are not incidental verbosity. The induction hypothesis must remain
quantified over the region, and the derivative assumption for $G_n$ is needed
only at $W(x)$ for $x\in U$. Replacing the latter by global differentiability
would simplify the interface by strengthening the theorem, not by explaining
the original theorem more naturally.

The existing proof is already well factored. Once this general theorem exists,
its applications may be shorter than any new narrative syntax. Clover's target
is the repeated assembly of locality and scope information in proofs of this
kind, not a claim that this particular theorem has no concise Lean formulation.

\subsection{Frey coefficients: one notation, two divisions}

The FLT source \code{S_FreyPackage_freyCurveInt_map.lean} proves that the
integral curve model maps to the rational one \cite{frey}. Its nontrivial
coefficient obligations compare
\[
 \left(\frac{b^p-1-a^p}{4}:\Z\right)\!\!:\Q
 \quad\hbox{with}\quad
 \frac{b^p-1-a^p}{4}:\Q,
\]
and the analogous expression with numerator $-a^pb^p$ and denominator $16$.
The source uses \code{Rat.intCast_div}, followed by the divisibility proofs
needed to justify the cast. The paper syntheses also exhibit the corresponding
\code{Int.cast_div} interface at their Lean/Mathlib pin.

The repetitive work has a precise shape: extract divisibility from the modular
assumptions, lift it through powers, normalize a residue modulo four, establish
exactness, and transport the resulting balance equation. Clover can package
that shape. It must not make integer division acquire field-division semantics
because their printed symbols look alike.

The same file also contains a short terminal \code{solution} that invokes the
helper. Comparing a short Clover call with the entire source file, including
other lemmas about invariants and valuations, would manufacture a misleading
compression ratio. The meaningful comparison is between equally equipped
interfaces at the same proof boundary.

\subsection{A finite tensor construction: properties of an entire family}

The FLT theorem in \code{S_Algebra_exists_algHom_equiv_pi.lean} constructs a
commutative $A$-algebra $W$ representing a finite family of algebra maps
\cite{tensor}. Its assumptions include that every factor $H_i$ is finite and
free as an $A$-module. Its conclusion includes both those module properties
for $W$ and a family of equivalences
\[
 e_T:\operatorname{AlgHom}_A(W,T)
       \simeq\prod_i\operatorname{AlgHom}_A(H_i,T)
\]
that is natural in $T$.

The code carries these data through induction on the finite index type,
transport along equivalences, the empty family, and the tensor-product step.
A property interface must preserve the dependency of $H_i$ on its index and
the naturality equation for the \emph{whole family} $e_T$. Independent tags
saying that the individual $e_T$ are equivalences would omit part of the
requested construction.

This example also tests the difference between evidence and data. The source
assumes \code{Finite} for the index type, not an author-supplied executable
enumeration. A mathematical existence proof may use a different interface
from a program required to compute a concrete tensor presentation.

\subsection{Leant: preserve the staging boundaries that already exist}

The directly inspected Leant source is pinned at
\code{3a40904be8a410d832d9ab6900b3d3e7b425eccb}. Its
\code{Verified} type has a nominal role and an explicit comment limiting its
meaning to acceptance by a supplied callback \cite{leant-verification}.
The Length contract module describes provider laws as passive assertions with
explicit argument roles. The adapter binds a checked candidate origin to a
specific Length problem and explicitly declines to treat a model-level failure
as authority for source-level behavioral refutation
\cite{leant-contract,leant-adapter}.

Both syntheses report a newer local Leant snapshot,
\code{823259f7e3c6}, whose behavioral path checks
an exact proposition about an exact candidate using Lean, and separately checks
its negation before reporting falsification. That is useful reported evidence;
it is not the same revision as the one directly fetched here
\cite{syn-a,syn-b}.

Clover should add a proof-retaining adapter at these boundaries, not weaken
Leant's existing distinctions between candidate identity, callback acceptance,
model evidence, and source semantics. Section~\ref{sec:synthesis} gives the
proposed contract.

\section{A small mathematical surface, not unrestricted prose}\label{sec:surface}

The full notation in this section is proposed syntax. Only the deliberately
smaller grammar in Section~\ref{sec:prototype} has an executable parser in this
package.

\subsection{Four ordinary authoring actions}

Clover distinguishes introducing data, assuming premises, establishing facts,
and constructing witnesses. A compact example is:
\begin{lstlisting}
fix f : A -> B
fix g : B -> A
assume inverse : left_inverse g f
know injective f by left_inverse
\end{lstlisting}
The first two lines introduce fixed objects. The third adds an explicit premise.
The fourth requests a proof of a consequence; it is not another assumption.
The proof transformation behind \code{by left_inverse} must produce a proof
of injectivity for that same $f$.

A refined binder is a convenient abbreviation:
\begin{lstlisting}
fix x : Real where x > 0
\end{lstlisting}
means an ordinary real variable together with a premise that it is positive.
By contrast,
\begin{lstlisting}
let y := expression where y > 0
\end{lstlisting}
creates an obligation to prove positivity of the specified expression. A
language that interprets both uses of \code{where} as assumptions would be
unsound as an authoring interface even if its final generated Lean declaration
were type-correct under the extra assumption.

Witness construction is visibly different:
\begin{lstlisting}
construct g : B -> A
  satisfying left_inverse g f
  using finite_inverse_construction
\end{lstlisting}
Here the author has authorized a choice of $g$, but not a change of $f$, $A$,
$B$, or the left-inverse specification. A successful construction returns data
and its proof. An unsuccessful attempt does not assert that no inverse exists.

\subsection{Operations and laws have different domains}

Clover uses \code{operation} for an interface with proof-implicit arguments:
\begin{lstlisting}
operation exact_quotient(n d : Int)
  requires nonzero : d != 0
  requires exact   : d divides n
  returns q : Int
  ensures d*q = n
\end{lstlisting}
Its ordinary-Lean meaning is a function accepting $n,d$, the required proofs,
and returning a quotient with a balance proof. A call generates those proof
arguments. This is an explicit partial-domain interface over Lean's total
integer-division primitive; it does not change the primitive itself.

A \code{law} instead states conditional behavior of a function that already
exists:
\begin{lstlisting}
law quotient_balance for integer_division
  given n d : Int
  requires d != 0 and d divides n
  ensures d * integer_division(n,d) = n
\end{lstlisting}
The carrier function remains defined at every pair of integers. Only the
advertised theorem is conditional. The two interfaces should not be conflated
by a single ambiguous arrow notation.

\subsection{Mathematical methods are typed transformations}

A method declaration associates a source operation with a theorem-backed
transformation of a proof problem. Representative methods are:
\begin{lstlisting}
differentiate identity on U
transport equation through ring_map phi
restrict action to invariant P
compare all finite observations
induct n keeping x in U general
\end{lstlisting}
These are not magic words. Each selects an interface with required premises and
a prescribed relationship between input statements and output statements. For
example, differentiating a regional identity must produce local equality at the
point of differentiation; restricting an action requires invariance of the
predicate for the \emph{selected} action.

A \code{hint} may guide proof search without constraining the accepted route.
A \code{by method} step requires the corresponding derivation node. A
\code{search} step delegates the route completely within the declared policy.
This makes the distinction visible without pretending that software can infer
a human author's intended proof method from an arbitrary proof term.

\subsection{A readable proof has a checked structural spine}

The proposed autonomous-derivative proof can be written as follows. Named
hypotheses retain the exact meanings of \eqref{eq:ode}--\eqref{eq:grec}.
\begin{lstlisting}
prove for n : Nat, on U:
  D^n W = G(n) compose W
by induction n keeping the point in U general:
  zero:
    use G_zero
  next n:
    differentiate IH on U
      using chain_rule from W_derivative and G_derivative(n)
    identify the result using G_step(n)
\end{lstlisting}
The mathematical spine is induction, local differentiation, and recurrence.
The expanded view contains the exact neighborhood fact, chain-rule application,
and rewrites. If the author removes openness, the system must not silently
accept the same method because its final target happens to have another proof.
It should identify the locality premise this method can no longer discharge.

Clover's prose view is generated from this structured proof document. Arbitrary
comments may still be written, but a displayed formula presented as an asserted
step must have its own evidence link. The final theorem's proof is not enough
to certify every sentence that happens to appear beside it.

\section{The author-facing type system}\label{sec:types}

\subsection{Carrier types, property interfaces, and evidence}

Let $\Gamma$ be an ordinary Lean local context, including selected structure
arguments. Let $E$ be an index into proof terms available in that context.
The index is not a second logical context: every accepted entry must have a
Lean type and proof in $\Gamma$.

A closed Clover judgment has the form
\[
 \Gamma;E\vdash t\Rightarrow A\;[\Phi],
\]
where $t:A$ is a fixed term and $\Phi$ is a finite interface of propositions
with retained proofs. A proposition in $\Phi$ may concern $t$ and other
objects; an injectivity fact concerns a function, while an inverse-pair fact
concerns two functions. A property interface is therefore not restricted to
unary labels such as \emph{positive} or \emph{continuous}.

Its interpretation is deliberately simple:
\[
 \Gamma\vdash t:A,
 \qquad
 \Gamma\vdash p_i:\phi_i \quad(\phi_i\in\Phi).
\]
Adding a property does not rebind $t$ to a new subtype. At an exported
interface, the same information may be packaged as
\[
 \CRef(A,P):=\{x:A\mid P(x)\},
\]
or as a named structure with data and proof fields. Lean already has these
representations; Clover makes their use-site requirements systematic
\cite{lean-subtype}.

There are four elementary rules. An accepted proof introduces a property;
a proved implication weakens or transforms an interface; properties about
one fixed configuration may be conjoined; and a checked equality may transport
a property. None of these rules lets a proposition become data merely because
it has a convenient English name.

For example, a proof that $a\ne0$ does not construct an inverse of $a$ in a
ring. In $\Z$, the element $2$ is nonzero and regular but not a unit. In
$\Z/4\Z$, the element represented by $2$ is nonzero but not regular. Clover
keeps the predicates \emph{nonzero}, \emph{regular}, and \emph{unit} distinct,
with conversions supplied only by theorems valid for the selected structure.

\subsection{Checking rather than guessing a stronger type}

Clover should not infer a mythical strongest refinement of an arbitrary term.
There are infinitely many true properties of even a simple numeral, and no
single finite interface captures all useful mathematics about it.

Instead, the system is bidirectional. In synthesis mode it obtains an
underlying type and a small interface exposed by the term's constructor,
provider specification, or known local facts. In checking mode an operation
requests a particular property interface. The obligation engine tries to
produce the missing proofs, using the rules allowed at that use site.

Thus an author may request
\[
 f:A\to B\quad[\mathsf{Injective}(f)],
\]
without requiring every function-valued expression to carry a permanently
materialized injectivity field. The semantic subsumption step is an explicit
proof-producing operation in elaboration, not an extra case in kernel
conversion.

This approach resembles refinement-type inference in its use of selected
qualifiers and implication obligations, but its predicates need not be confined
to a decidable SMT fragment. Outside a chosen automatic fragment, they remain
ordinary proof obligations. Liquid Types and refinement-type inference provide
relevant design precedents, not a ready-made decision procedure for arbitrary
Lean mathematics \cite{liquid,refinement-tutorial}.

\subsection{Pure behavioral contracts}

For a fixed total function $f:A\to B$, define
\[
 \Beh(f;P,Q):=\forall x:A,\ P(x)\to Q(x,f(x)).
\]
The postcondition is relational: it can refer to the actual input and output.
A sorting contract is not merely that the output is sorted. It includes
permutation of the input:
\[
 \forall xs,\quad
 \mathsf{Sorted}_{\le}(f(xs))\ \land\
 \mathsf{Perm}(xs,f(xs)).
\]
Stability, a comparison-cost bound, or preservation of a separate invariant
are additional requirements, not consequences of the bare function type.
A cost claim needs a specified operational cost model; mathematical extensional
correctness does not imply it.

Variance is proof-directed. Suppose $P'(x)\Rightarrow P(x)$ and
$Q(x,y)\Rightarrow Q'(x,y)$. A contract from $P$ to $Q$ yields a contract
from $P'$ to $Q'$. The precondition becomes more restrictive and the guarantee
weaker. There is no scalar ranking in which every possible contract is simply
``stronger'' or ``weaker'' than every other one.

Composition preserves the intermediate value. If
\[
 \Beh(f;P,Q),\qquad \Beh(g;R,S),
\]
then a useful composition rule requires
\[
 \forall x,y,\ P(x)\land Q(x,y)\Rightarrow R(y).
\]
It yields, for example,
\[
 \forall x,\ P(x)\Rightarrow
 Q(x,f(x))\land S(f(x),g(f(x))).
\]
A final output relation $T(x,g(f(x)))$ requires one more proved implication.
Replacing the actual $f(x)$ by an unrelated existential witness would lose the
connection that justifies the second call.

Higher-order and relational properties use the same discipline. Monotonicity,
Lipschitz continuity, naturality, equivariance, and commutation of two maps are
propositions about multiple evaluations or whole families. They are not
approximated by a collection of unrelated unary output predicates.

\subsection{Refined domains are not guarded total functions}

A function
\[
 f:\{x:A\mid P(x)\}\to B
\]
need not extend to $A\to B$. There may be no suitable output outside the
refined domain, or the desired extension may require a choice of default or
a separate construction. A guarded law about an existing total function does
not have this difficulty because the data already exist everywhere.

Likewise, a proof of $\exists y, Q(y)$ is not generally an executable procedure
returning such a $y$. Lean distinguishes proposition-valued existence from
computational data; classical choice can provide a mathematical selection, but
that is a different construction policy \cite{lean-quantifiers}. Clover's
\code{construct} node records whether its result is computational or invokes
an admitted noncomputable selection.

\subsection{Existing Lean mechanisms are part of the baseline}

Lean already supports automatic parameters whose associated tactic scripts
supply missing arguments. It also supports typeclass inference and ordinary
bundled properties \cite{lean-application}. Consequently, the mere syntax
\code{requires} is not evidence of a new capability. A first implementation
could lower it to automatic proof arguments backed by a shared resolver.

The proposed contribution beyond that baseline is the coordinated handling
of quantified scope, a shared and inspectable obligation graph, role-checked
method interfaces, controlled rule generation, and diagnostics that preserve
the mathematical meaning of the missing premise. Whether that coordination
justifies separate syntax is an empirical question.

\section{Unfinished reasoning as a dependent obligation telescope}\label{sec:obligations}

\subsection{A pending result is a function awaiting evidence}

An unfinished elaboration judgment has the form
\[
 \Gamma;E\vdash e\Rightarrow A[\Phi]\;!\;\Delta,
\]
where
\[
 \Delta=(h_1:P_1,\ h_2:P_2(h_1),\ldots,
             h_m:P_m(h_1,\ldots,h_{m-1}))
\]
is a well-formed telescope of proof obligations. In the core fragment, each
binder in $\Delta$ is proposition-valued. Authorized data constructions are
separate nodes that extend $\Gamma$ with a chosen witness and the evidence
that accompanies it.

The meaning is that a skeleton $e:A$ and the advertised property proofs exist
under $\Gamma,\Delta$. Before the obligations are discharged, the exported
object is at most the abstraction
\[
 \lambda h_1\ldots h_m.\,e : \Pi\Delta, A,
\]
not an unconditional term of $A$. This is the precise sense in which unfinished
proof work behaves like an \emph{obligation effect}. It is not a new runtime
effect and does not let the program read or manufacture evidence.

A pending expression may itself require proof arguments to be well typed.
Its placeholder is therefore a suspended elaboration object, not an ordinary
closed anchor that may already be used in arbitrary accepted facts. This
restriction prevents a superficially convenient implementation from placing
properties of ill-typed or incompletely specified data into the evidence index.

\subsection{Sequential composition is dependent substitution}

Suppose the first step produces a skeleton $y=e(x)$ and obligations $\Delta$.
The next method asks for a property $R(y)$ and produces obligations $\Theta(y)$.
Composition forms the telescope
\[
 \Delta\ ;\ \Theta(e(x)),
\]
with the actual intermediate substituted into every dependent requirement.
If a proof in $\Delta$ is solved, it is substituted into later obligations;
it is not simply deleted from a flat set of strings.

For exact quotient transport, the sequence is naturally ordered. First fix
the integer quotient expression. Then prove its balance equation using the
divisibility hypothesis. Map that equation to the target ring. Finally, if
the requested conclusion contains division in that ring, supply a unit or
field-nonzero premise for the mapped denominator. The final requirement is
about the mapped denominator, not about an unrelated integer proposition with
a similar name.

\subsection{Discharging an obligation does not authorize a cycle}

Assume a method gives a conditional theorem $P\to Q$. The implementation may
retain this theorem while $P$ remains open. It may not place $Q$ among the
unconditional facts and then use $Q\to P$ to solve the original guard. The two
implications alone establish neither proposition.

Clover therefore distinguishes a \emph{rule dependency graph}, which can have
cycles, from an \emph{accepted proof graph}, which is built from already
available proofs and has no unsupported circular dependency. Induction and
well-founded recursion are not forbidden: they are explicit logical rules
with their own checked premises, not cycles in ordinary fact propagation.

A discharge operation must either supply a proof in the appropriate earlier
context or retain the obligation. A branch-local proof is valid only under
that branch's telescope. Leaving the branch exports a conditional theorem or
a proof formed by case elimination; it does not export the local assumption
as a fact in the parent scope.

\subsection{The core translation theorem}

\begin{theorem}[Conditional preservation for the specified core]\label{thm:preservation}
Assume that each registered method is accompanied by an ordinary Lean theorem
of its declared type, that substitutions are well typed, and that object and
structure arguments have been fixed as specified. A derivation of
\[
 \Gamma;E\vdash e\Rightarrow A[\Phi]\;!\;\Delta
\]
using the core introduction, application, refinement, transport, scope, and
obligation-composition rules translates to an ordinary Lean skeleton of type
$A$ in $\Gamma,\Delta$, with proofs of the propositions in $\Phi$ there.
A well-typed closing substitution for $\Delta$ yields an ordinary Lean term
of the advertised type in $\Gamma$.
\end{theorem}

\begin{proof}
Proceed by induction on the derivation. Variable and premise cases use the
corresponding declarations in the context. A refinement-introduction case
pairs the unchanged term with an available proof at the elaboration level;
packaging at an interface uses subtype or structure introduction. Consequence
applies the retained implication theorem. A method application instantiates
its Lean theorem and applies it to the translated premises, introducing a
proof binder for every still-open premise. Sequential composition uses the
ordinary substitution lemma for dependent function types. Equality transport
uses equality elimination with the actual predicate and indices. Branch entry
extends the context; branch exit abstracts its assumptions or applies the
specified case-elimination theorem. These operations produce terms in the
claimed telescope. Applying a closing substitution to the resulting terms
preserves their types and gives the final claim.
\end{proof}

This is a mathematical theorem about the stated rules. It is not a machine
verification of the proposed parser, rule metadata validator, renderer, or
scheduler. In particular, the theorem assumes the very object correspondence
that a production implementation must establish. The more informative
implementation obligations are the correspondence and state-transition tests
in Sections~\ref{sec:meaning}, \ref{sec:transport}, and \ref{sec:prototype}.

\subsection{An obligation frontier that explains the next decision}

A frontier entry should contain the exact proposition, its local context,
its source span, the operation consuming it, and available sufficient routes.
For example:
\begin{lstlisting}
Step: transport q = n/d from Int to K
Open requirement: IsUnit (phi d)
Already available: d divides n; d != 0 in Int
Reason: this route rewrites the mapped balance as division in K
Alternative: keep the balance equation (phi d)*(phi q) = phi n
\end{lstlisting}
The last line describes an alternative \emph{result}, not a silently weakened
answer. Switching to that result changes the request and requires an author
choice. Likewise, a premise listed by one route is not announced as necessary
for every proof of the original theorem.

The first implementation need not find the smallest frontier or the weakest
sufficient assumptions. Those optimization problems can be more expensive
than finding a proof. It should report one or a few traceable routes and keep
the original target fixed.

\section{Quantified evidence and mathematical scope}\label{sec:quantified}

\subsection{Store a theorem scheme, not a cloud of instances}

An evidence index must retain facts such as
\[
 h:\forall n\in\N\;\forall x\in U,\ P(n,x)
\]
as quantified schemes. Storing only the instances already used loses the
information needed for later reasoning under a new arbitrary point or index.
Conversely, eagerly enumerating all instances is impossible even for the
simplest infinite domain.

The index stores the binder telescope and its dependencies. A use site selects
an instance by matching the requested proposition, supplying data arguments
from the current scope, and proving its guards. The universal proof remains
available unchanged. This is ordinary universal elimination with an operational
index, not a new axiom about schematic formulas.

For the autonomous derivative proof, the induction step begins with
\[
 \mathrm{IH}:\forall y\in U,\ D^nW(y)=G_n(W(y)).
\]
At an arbitrary $x\in U$, openness gives a neighborhood contained in $U$.
Specializing IH \emph{throughout that neighborhood} supplies the eventual
equality needed for differentiation. Specializing IH to $x$ first and forgetting
its quantified form cannot justify the same step.

\subsection{Region syntax is a binder discipline}

The proposed block
\begin{lstlisting}
on U:
  know f = g
\end{lstlisting}
means a checked universally quantified regional equality, not a rebinding of
$f$ and $g$ to functions on a different domain. A local block
\begin{lstlisting}
at x in U:
  know f(x) = g(x)
\end{lstlisting}
is weaker unless the surrounding proof establishes that $x$ is arbitrary and
exports the appropriate quantifier. The document model records that distinction
before the pretty printer renders either block.

The same principle applies to dependent families. In a proof of
$\forall n,\forall x\in U_n,P(n,x)$, the region depends on the index. An
induction method must state what it generalizes and what it fixes; it cannot
borrow an induction hypothesis over $U_n$ as if it covered $U_{n+1}$ without
an inclusion or transport theorem.

\subsection{Quantifier order is not decorative}

A particularly sharp example comes from measure theory. On $[0,1]$ with
Lebesgue measure, define $f_t(x)=\mathbf 1_{\{t\}}(x)$ and $g_t(x)=0$. For
every fixed $t$, the two functions are equal almost everywhere. But there is
no $x\in[0,1]$ at which they are equal for every $t$, because $f_x(x)=1$.
Thus
\[
 \forall t,\quad f_t\AEeq g_t
 \quad\not\Longrightarrow\quad
 \text{for almost every }x,\ \forall t,\ f_t(x)=g_t(x).
\]
For a countable parameter family, the implication does hold: intersect the
countably many full-measure sets. This gives a precise contract for a
\code{collect almost_everywhere} method: it needs countability or another
proved uniform-null-set condition.

Clover should represent \emph{for each parameter} and \emph{almost everywhere}
as different binders, not as unordered tags in a property row. This example
also shows why the generic term ``local'' is too coarse. Neighborhood equality,
regional equality, and almost-everywhere equality have different consumers and
different quantifier laws.

\subsection{Almost-everywhere consumers are relation-specific}

For a fixed measure $\mu$, an a.e. equality $f=_{\mu\text{-a.e.}}g$ supports
an integral congruence through the appropriate theorem for the chosen integral.
It does not support evaluation at a chosen point. If the interface also claims
integrability or a finite integral, those properties must be carried through
the corresponding integrability theorems; totalized notation alone does not
supply them.

A consumer registration records the measure argument and the relation being
used. It should distinguish a theorem that respects a.e. equality from a
stronger theorem that extracts pointwise equality under continuity and support
hypotheses. The latter is a separate mathematical result, not a coercion that
can be inserted whenever point evaluation would be convenient.

\subsection{Strategic steps: symmetry and without loss of generality}

A language aimed at mathematical reasoning must support more than local fact
inference. The phrase ``without loss of generality'' is a particularly useful
test because it deliberately changes the configuration while preserving the
original proof obligation.

Let $X$ be a configuration type, $P:X\to\Prop$ its admissibility condition,
$Q:X\to\Prop$ the goal, and $N:X\to\Prop$ a preferred normal position.
A normalization method supplies $r:X\to X$ and proofs of
\begin{align*}
 P(x)&\Rightarrow P(r(x)),\\
 P(x)&\Rightarrow N(r(x)),\\
 P(x)&\Rightarrow\bigl(Q(r(x))\Rightarrow Q(x)\bigr).
\end{align*}
A proof of $\forall y,\ P(y)\land N(y)\Rightarrow Q(y)$ then proves the
original $\forall x,\ P(x)\Rightarrow Q(x)$: apply it to $r(x)$ and use the
last implication. These are the exact premises of the method, not a license
to replace $x$ by an equal object.

For a symmetric statement about two elements of a linear order, $r$ can keep
$(x,y)$ when $x\le y$ and swap it otherwise. The method needs preservation
of the hypotheses under swapping and the implication transporting the goal
back. A proposed Clover step is:
\begin{lstlisting}
without_loss x <= y
  using swap(x,y)
  preserving assumptions
  transporting goal by symmetry
\end{lstlisting}
The continuation works on the selected normalized configuration, while the
outer declaration still proves the original statement. If the assumptions
are not symmetric, the method must report that missing preservation proof;
it may not copy them unchanged to the swapped variables.

This rule also explains the interaction between meaning preservation and
legitimate data choice. The original target is sealed, but an explicitly
selected proof transformation may construct new auxiliary data and prove a
reduction back to that target. Freezing meaning is not a prohibition on normal
mathematical argument; it is a prohibition on changing the claim without such
a reduction. The general strategic rule is specified here, not implemented
by the small endomorphism parser.

\section{Fixing meaning without freezing legitimate constructions}\label{sec:meaning}

\subsection{Three kinds of holes}

Clover classifies unfinished information by its semantic role rather than by
whether it happens to be represented internally as a metavariable.

\begin{center}\small
\begin{tabularx}{\textwidth}{@{}P{.19\textwidth}Y Y@{}}
\toprule
Role & Examples & Permitted resolution\\\midrule
Meaning-bearing & Carrier type, topology, measure, coefficient ring, branch, proposition parameter & Ordinary statement elaboration; unresolved ambiguity is reported or becomes an explicit rigid parameter.\\
Authorized data & A requested inverse, basis, factorization, algorithm, or representing object & A construction provider may choose a witness satisfying the frozen specification.\\
Proof & Nonzeroness, divisibility, continuity, a contract proof & Any admitted proof-producing provider may discharge the exact obligation.\\
\bottomrule
\end{tabularx}
\end{center}

The distinction is necessary because a proposition can be made easy by choosing
its missing data. A solver asked to prove $x\ge0$ must not set an unresolved
meaning-bearing $x$ to $0$ unless the source requested construction of \emph{some}
nonnegative real. On the other hand, a request to construct a polynomial
factorization may legitimately leave the factor coefficients as data holes.

\subsection{The proposed sealing protocol}

Statement elaboration first resolves notation and selected structures and
produces a typed statement skeleton. The sealer then computes a dependency
closure of the remaining metavariables: types of holes, universe constraints,
instance arguments, and data occurring in the target all matter. Looking only
for metavariables printed in the conclusion is insufficient.

Every remaining meaning-bearing component is either resolved before proof
search, made an explicit rigid parameter of an authorized generalized request,
or reported as an ambiguity. Making a parameter rigid must preserve the
intended statement; it cannot silently generalize a source that asked for a
particular object. The baseline implementation should prefer a diagnostic when
that intent is not determined.

Proof and construction providers run in copied elaboration state. Their
results are checked against the sealed target and the declared hole roles.
The commit step examines assignments and rejects unauthorized changes to
meaning-bearing variables or their dependency closure. A proof provider may
create internal temporary metavariables, but every temporary needed by the
accepted result must be solved or explicitly returned as an obligation.

Lean's elaboration APIs distinguish postponed and synthetic metavariables,
but their internal categories should not be mistaken for this entire semantic
policy. In particular, using an opaque proof metavariable does not by itself
freeze every data or universe variable reachable from its type
\cite{lean-synthetic,lean-elaborators}.

\subsection{Preserve the chosen structure, not just its carrier}

A topology, order, multiplication, scalar action, or measurable structure is
mathematical data. Two such structures on one carrier need not be equal. A
proof that a fixed function is continuous under one topology does not apply
under another topology merely because the function and carrier print the same.

A Clover lexical structure context records the actual elaborated dictionaries
used in the block. The context is not an instruction to bypass Lean's instance
system globally. It supplies explicit choices at selected interfaces and
creates local bridges only where a library theorem expects an instance.

This policy may make generated instance-management preambles easier to
understand, but no such reduction was measured here. The inspected Frey and
tensor files are not evidence that every FLT file has a large avoidable
preamble. Replacing one actual preamble is a separate experiment.

\subsection{A precise but limited noninterference property}

At the specification level, proof search should satisfy: if two accepted runs
start from the same sealed statement and differ only in proof-provider choices,
the resulting declarations have the same advertised proposition, up to the
explicitly admitted equivalences of the sealing protocol. Different witnesses
are allowed only in authorized construction positions.

The easiest first implementation enforces a stronger operational condition:
no assignments to the protected metavariable closure, exact checking of the
retained target, and no migration of local identities across rebuilt contexts.
This can reject harmless changes, but it gives an implementable initial
boundary. A later implementation may accept more cases through checked context
or statement adapters; a matching hash alone is never such an adapter.

\section{A finite inference engine with specified rule generation}\label{sec:grounding}

\subsection{What a registry entry contains}

A registered rule is an ordinary theorem plus elaboration metadata. Its
interface identifies the conclusion pattern, data parameters, premise
parameters, structure parameters already fixed by the context, and permitted
trigger positions. The metadata validator must compare those roles with the
actual elaborated theorem type. A misspelled role annotation is an error, not
permission to reinterpret the theorem.

The first automatic fragment uses Horn-shaped theorem interfaces:
\[
 \forall\vec a,\quad P_1(\vec a)\to\cdots\to P_k(\vec a)\to Q(\vec a).
\]
All data in the instantiated propositions are fixed terms. The automatic
fragment excludes general synthesis of new data from a proof premise,
unrestricted higher-order unification, and dependent argument patterns whose
meaning cannot be determined before their proofs are found. Those cases remain
available through explicitly typed method applications and ordinary Lean
elaboration; they are not promised to lie in the finite closure engine.

For example, the injectivity rule
\[
 \mathsf{LeftInverse}(g,f)\to\mathsf{Injective}(f)
\]
has a data parameter $g$ that occurs only in its premise. Searching for an
inverse of $f$ is not the same operation as instantiating an existing theorem
with one of the functions already present. The finite engine does only the
latter. A new inverse is an authorized construction request.

\subsection{The term arena}

For each supported data type, construct a finite arena from the elaborated
terms occurring in the requested statement, visible assumptions, selected
method arguments, and retained evidence relevant to the block. Include their
well-typed subterms. Keep binder identity, universe instantiations, and
structure arguments; do not identify terms because their pretty-printed
strings agree.

The initial arena does not close under arbitrary function composition,
polynomial formation, or constructor application. If $f$ and $g$ occur, that
alone does not generate $g\circ f$, $f\circ g$, and all their iterates. A
composition explicitly present in the target is already in the arena. A
registered constructor may be used in a later, separately bounded widening
phase or under an explicit \code{consider} request, but it must not create an
unbounded hidden term universe during closure.

This restriction sacrifices some automatic proofs. It is also what makes the
inference boundary intelligible: the engine is assembling consequences about
a known set of mathematical objects, not searching simultaneously for every
possible auxiliary object.

\subsection{Goal-triggered grounding}

The following procedure gives a complete finite grounding policy for the
chosen arena and rule profile. ``Match'' below means a terminating, selected
pattern-matching procedure, not arbitrary semantic unification.

\begin{lstlisting}
work := [requested proposition]
seen := empty set of propositions
instances := empty set of ground theorem applications

while work is not empty:
  pop a proposition Q
  if Q is already in seen: continue
  add Q to seen
  for each enabled rule whose conclusion matches Q:
    bind the parameters fixed by that conclusion match
    enumerate remaining data parameters over their finite arenas
    retain only well-typed instances using arena terms
    record the exact theorem instance
    add its premise propositions to work

compute the forward closure of available proved facts
using these recorded instances
\end{lstlisting}

A candidate limit, term-size limit, or unsupported pattern can stop this
procedure. Such a stop is a resource or fragment outcome. It must not be
reported as completed finite non-derivability. In the reference implementation,
exceeding the instance-attempt limit returns \code{resource_limit} and never
returns a negative mathematical verdict.

A parameter that appears only in premises is enumerated from the arena. A
parameter whose type depends on earlier data parameters is drawn only after
those parameters have been instantiated and its type has been checked. If no
finite arena for that dependent type is available, this rule application is
outside the automatic fragment. Structure arguments are taken from the sealed
context rather than guessed from a global collection of alternative structures.

The production implementation should attempt exact evidence and cheap
single-lemma matches before materializing the whole relevant closure. It can
return immediately after a checked proof is found. Exhaustive grounding is
needed only for the stated finite-fragment closure result, not as a condition
for accepting every successful proof.

\subsection{Closure, proof graphs, and a finite guarantee}

Let $F_0$ be the available proved atoms and let $\mathcal G$ be the finite
ground-rule set. Define
\[
 F_{n+1}=F_n\cup
 \{Q\mid (P_1,\ldots,P_k\Rightarrow Q)\in\mathcal G,
                    \ P_1,\ldots,P_k\in F_n\}.
\]
Whenever a new atom is inserted, retain a proof node containing the exact rule
instance and references to proofs already available for its premises.

\begin{theorem}[Finite-fragment closure]\label{thm:closure}
For a fixed finite arena, a finite rule profile with terminating pattern
matching, and completed grounding, the closure process terminates. Its final
set is the least rule-closed superset of $F_0$. Any fair schedule reaches the
same set of atoms. If the seed proofs and instantiated rule theorems are valid,
every inserted atom has a valid proof. The selected proof, dependency support,
and running cost need not be schedule-independent.
\end{theorem}

\begin{proof}
Only finitely many instantiated atoms occur in $F_0$ and $\mathcal G$. Each
strict change adds an atom, so there can be only finitely many changes. The
limit is closed because every rule whose premises are present eventually fires.
Every closed superset of $F_0$ contains each $F_n$, by induction, establishing
leastness. A fair schedule cannot omit a consequence at the least fixed point,
again by induction on the stage at which that consequence first appears.
Each proof node applies an actual instantiated theorem to already justified
premises, so induction on insertion order proves validity. Different schedules
may choose different justifications for the same atom, which explains the last
qualification.
\end{proof}

This is relative completeness for a \emph{selected finite rule system}. It says
nothing about all true propositions of Lean, all possible qualifier choices,
or all auxiliary constructions a mathematician might invent.

With at most $M$ arena terms per relevant type and at most $v$ data parameters
per rule, a simple upper bound on possible substitutions is $rM^v$ for $r$
rules, before accounting for dependent typing and pattern costs. This is a
finite bound, not an attractive complexity guarantee. With a premise-indexed
work queue and a residual-premise counter per ground rule, closure bookkeeping
can be linear in the number of atoms and ground premise occurrences. Grounding,
Lean type checking, normalization, proof retention, and rendering remain
separate costs.

\subsection{Profiles are interfaces, not global mathematical omniscience}

A profile should name a small, coherent set of theorem interfaces: exact integer
quotients, regional differentiation, finite affine lengths, or injectivity of
compositions. It can import other profiles explicitly. Automatically enabling
every theorem that resembles a property rule would recreate an unpredictable
second elaborator.

The companion supplies a concrete warning. For one composition problem, the
seven-rule registry generated 211 ground instances and retained a four-node
proof. Restricting it to the two rules actually needed reduced that count to
16 while retaining a four-node proof. These are counts in the restricted Python
model, not measurements of Lean inference. They nevertheless demonstrate that
``finite and demand-driven'' does not imply ``cheap'': equality symmetry and
transitivity can create a large relevant-looking search region even when the
final proof never uses equality.

A sensible production policy is therefore staged: direct lookup; registered
single-step rules; a small method-specific composition profile; explicit
widening; finally broader external search. Each stage must retain its proof
and outcome classification. No unsuccessful stage acquires authority to negate
the target.

\section{Transport, rewriting, and the lifetime of evidence}\label{sec:transport}

\subsection{Three distinct ways to reuse a fact}

Suppose an index contains $p:P(t)$ and a consumer asks for a fact about $u$.
There are three common legal routes.

If $t$ and $u$ are definitionally equal in the fixed Lean context, ordinary
type checking can reuse $p$. If a checked equality $e:t=u$ is available,
equality elimination transports $p$ to $P(u)$. If only a weaker relation
$R(t,u)$ is known, reuse needs a consumer theorem of an appropriate shape, such
as
\[
 R(t,u)\to P(t)\to Q(u).
\]
The third route does not assert $t=u$ and must not populate the equality index
as if it did.

The same applies when a simplifier rewrites propositions rather than their
objects. From $p:P(t)$ and a proved equivalence $P(t)\leftrightarrow Q(u)$,
the implication gives $Q(u)$ without proving $t=u$. The index may keep its
original anchor and insert this bridge at the use site; it need not continually
rewrite every stored fact into one global normal form.

\subsection{Dependent indices make transport a typed operation}

For $v:\mathsf{Vector}(A,n)$ and an equality $n=m$, a transported vector has
a different indexed type. A property about the original vector is not reused
by deleting its length index or matching a textual name. The adapter must
carry the index equality, the transported value, and the predicate transport
needed by the consumer.

Proof irrelevance can eliminate distinctions between proofs of a fixed
proposition; it does not identify distinct data indices or structure arguments
\cite{lean-propositions}. The implementation should initially support a few
explicit transport shapes, keeping ordinary Lean as an escape hatch for
more complicated dependent equalities.

A normalization policy must likewise distinguish an exact equality from a
chosen isomorphism. Isomorphic vector spaces or representing algebras are not
interchangeable data in every dependent expression. A library can supply
coherent transport along an isomorphism, but the chosen map and its inverse
remain part of that operation's meaning.

\subsection{Snapshot-local indexing is the initial lifetime rule}

Each evidence index belongs to an elaboration snapshot. Within it, keys use
actual expressions and local identities. After an upstream edit rebuilds a
declaration, the initial implementation reconstructs the index from the new
Lean context rather than reusing old local identifiers.

More ambitious migration requires a typed context embedding. If
\[
 \Gamma=(x_1:A_1,\ldots,x_n:A_n)
\]
and the new context is $\Gamma'$, a migration map must provide terms
$\theta(x_i)$ whose types are the corresponding substitutions of $A_i$,
in dependency order. Replaying a proof then means substituting $\theta$ into
that proof and checking the result. A hash match or a matching source name is
useful for locating a candidate migration, not for proving it correct.

This conservative rebuilding policy deliberately misses some reusable work.
The restricted companion is even coarser: its snapshot key includes the entire
source text, so even an unrelated edit invalidates old request identities.
That is a safe prototype boundary, not the proposed final incremental editor
architecture.

\subsection{A framing law and its limits}

Within a fixed signature and fixed data context, adding valid evidence cannot
invalidate an already accepted proof. It can make new obligations solvable.
For a fixed finite arena and rule profile, the closure operator is monotone in
its seed facts. This is the evidence analogue of a framing property.

The qualification matters. Adding an import can change notation resolution or
available instance choices during a \emph{new} elaboration. Changing a method
policy can make an old derivation unacceptable even when its theorem remains
true. Changing a measure changes the proposition. These are not counterexamples
to monotonicity because the sealed signature or acceptance policy has changed.

Clover should make that distinction visible in repair diagnostics: an old proof
may fail because a fact disappeared, because the statement changed, or because
the permitted method changed. A single undifferentiated ``automation failed''
message conceals the very information needed to repair the proof.

\section{Behavioral synthesis and realizable observation types}\label{sec:synthesis}

\subsection{The synthesis request is a dependent result type}

For a requested program type $T$ and fixed specification $S:T\to\Prop$, the
accepted result is conceptually
\[
 \{f:T\mid S(f)\}.
\]
A producer may search by types, examples, symbolic constraints, or language-model
suggestions. The acceptance interface does not care which heuristic found $f$;
it requires evidence of the \emph{same} $S(f)$ for the \emph{same} candidate.

A proposed Leant adapter retains three linked objects: the selected source edit
and its elaborated term; the sealed specification with its instantiated types,
universes, and structures; and a Lean proof of that specification or a checked
certificate bridge producing one. Candidate identity and theorem correspondence
are not reconstructed from a provider name or a pretty-printed expression.

The existing nominal callback receipt remains useful for staging. It should
not be repurposed as the new evidence type. The additional receipt must state
which target was checked, which proof or certificate is retained, and which
axiom and execution policy was used. Replaying a different textual variant
requires a new check of the edited source.

\subsection{Soundness, completeness, and refutation are different contracts}

Let $D$ be an admissible concrete input domain, let $o:D\to U$ be an observation,
and let an abstract evaluator describe the candidate on observations. Four
questions must be answered independently.

\textbf{Positive soundness:} does acceptance of the abstract certificate imply
$S(f)$ for all admissible inputs? A sound but incomplete method is already
useful here.

\textbf{Decision completeness:} does the method accept every true specification
in its stated program grammar and input domain? This is a stronger converse,
not a prerequisite for accepting sound proofs.

\textbf{Negative transfer and realization:} does a reported abstract failure
come from an admissible concrete input, and would the source specification
imply the abstract property at that input? Both are needed for refutation.

\textbf{Observation coverage:} why do the available observations determine the
property being claimed? A theorem over all finite prefixes is not the same
thing as having checked several prefixes.

Clover represents these as different proof-bearing interfaces. A provider with
only positive soundness can prove accepted cases. It may not label its failure
as a counterexample. A provider with a concrete source counterexample may
refute that candidate without possessing a complete decision procedure for the
whole grammar.

\subsection{A more general finite test basis for affine observations}

The usual unrestricted affine-length criterion compares coefficients. That
criterion is too strong for some restricted domains and can produce unrealizable
negative witnesses. The right object is the affine geometry of the
\emph{realizable observation image}.

The following elementary theorem provides a precise reusable interface. Its
value here is not a claim of new linear algebra; it is the way its hypotheses
can be packaged as behavioral type information for synthesis.

\begin{definition}[Realizable affine test basis]
Let $K$ be a field, $D$ a concrete admissible domain, and $o:D\to K^m$ an
observation map. A realizable affine test basis consists of a base point $b$,
directions $v_1,\ldots,v_r\in K^m$, and concrete inputs
$x_0,x_1,\ldots,x_r\in D$ such that
\[
 o(x_0)=b,\qquad o(x_i)=b+v_i\quad(1\le i\le r),
\]
together with the coverage property
\[
 \forall x\in D,\quad
 o(x)-b\in\operatorname{span}_K\{v_1,\ldots,v_r\}.
\]
The directions need not be linearly independent.
\end{definition}

\begin{theorem}[Affine specification from a realizable test basis]\label{thm:basis}
Suppose $o$ has a realizable affine test basis. For affine functions
$p,q:K^m\to K$,
\[
 \bigl(\forall x\in D,\ p(o(x))=q(o(x))\bigr)
 \quad\Longleftrightarrow\quad
 \bigwedge_{i=0}^{r} p(o(x_i))=q(o(x_i)).
\]
If one of these finite equalities fails, the corresponding $x_i$ is an
admissible concrete counterexample to the universal equality.
\end{theorem}

\begin{proof}
Write $p(u)-q(u)=c+\ell(u)$ for a linear functional $\ell$. Equality at $x_0$
gives $c+\ell(b)=0$. Equality at $x_i$ then implies
\[
 0=c+\ell(b+v_i)=c+\ell(b)+\ell(v_i)=\ell(v_i).
\]
For arbitrary $x\in D$, coverage expresses $o(x)-b$ as a linear combination
of the directions, so $\ell(o(x)-b)=0$. Hence
\[
 p(o(x))-q(o(x))=c+\ell(b)+\ell(o(x)-b)=0.
\]
The reverse implication specializes the universal claim to the supplied inputs.
A failing equality is already a failure at a named member of $D$, giving the
last assertion.
\end{proof}

The proof separates useful assumptions. Coverage alone, together with checked
values at the formal points $b,b+v_i$, suffices for a positive affine identity
on the image. Concrete realizers make those checks necessary for the universal
source property and turn failures into source witnesses. Full surjectivity of
$o$ onto $K^m$ is unnecessary. If $D$ is empty, the universal statement is
vacuously true and is handled separately; there is no required base-point
realizer.

\subsection{Correlated lists: three observations, only two directions}

Take $D=\mathsf{List}(A)\times\mathsf{List}(A)$, with a chosen inhabitant
$a\in A$, and observe
\[
 o(xs,ys)=(|xs\mathbin{+\!+}ys|,\ |xs|,\ |ys|)\in\Q^3.
\]
The image lies in the plane $n=p+q$. A realizable affine test basis is
\[
 b=(0,0,0),\qquad v_1=(1,1,0),\qquad v_2=(1,0,1),
\]
realized by $([],[])$, $([a],[])$, and $([],[a])$. Every observation is
$|xs|v_1+|ys|v_2$.

Consequently, two affine expressions in total length, left length, and right
length agree for all such list pairs exactly when they agree at these three
realizable points. Their coefficient vectors in three independent variables
need not be equal: $n$ and $p+q$ are the simplest example. An abstract vector
such as $(1,0,0)$ violates the observation relation and cannot refute the law.

For a difference
\[
 d(n,p,q)=c+an+bp+eq,
\]
the complete criterion is
\[
 c=0,\qquad a+b=0,\qquad a+e=0.
\]
The companion enumerates 625 small integer-coefficient forms, checks this
criterion against the finite test basis, and reconstructs actual list-pair
witnesses for every failed basis check. It also performs finite grid checks
for accepted forms. The universal theorem is the proof above; the finite run
is an implementation check, not its replacement.

\subsection{Empty element types and guards change the basis}

If $A$ is empty, both input lists must be empty. The observation image is just
$\{(0,0,0)\}$, so a basis has one base point and no directions. This validates
some universal source equations whose unrestricted affine coefficient vectors
differ. It also explains why an abstract positive-length failure cannot be
promoted to a source counterexample at \code{List Empty}.

A guard can have the same effect without an empty element type. On the domain
of lists of length exactly $100$, the length observation has image $\{100\}$.
The affine expressions $n$ and $100$ agree on that domain although they are
not identical affine functions on all natural lengths. A realizer requires
an actual list of length $100$; it must not be fabricated when the element
type is empty.

More generally, a length guard or correlation defines a constrained observation
image. Clover should ask for an adequacy package for that image rather than
reuse an unrestricted coefficient criterion by default. There need not be a
convenient computable basis for every guard. Failure to construct one leaves
the method unsupported, not the specification false.

\subsection{How this attaches to real candidate semantics}

The test-basis theorem does not itself connect Leant's term graph to affine
lengths. A complete adapter still needs an interpreter for the admitted
candidate grammar, a normalization theorem computing its affine observation,
and a correspondence theorem for the exact elaborated source candidate.
Provider laws must be Lean theorems about the actual provider terms, not passive
model assumptions promoted by name.

The sound path is therefore
\[
\begin{gathered}
 \text{exact candidate}\ \longrightarrow\
 \text{checked denotation bridge}\ \longrightarrow\
 \text{affine model}\\[3pt]
 \longrightarrow\ \text{image-aware certificate}\ \longrightarrow\
 \text{source contract proof}.
\end{gathered}
\]
For a negative verdict, add the checked source input realizer and evaluate or
prove the failure on that input. The companion's affine checker does not claim
to implement this full path.

A polymorphic-looking type is not by itself a behavior theorem. Lean's
reference gives a classical, noncomputable polymorphic list function that
inspects its element type, invalidating an unrestricted parametricity rule
\cite{lean-axioms}. A restricted synthesis grammar can support a separately
proved parametricity theorem; membership in that grammar then becomes part of
the evidence, not an inference from the printed function type.

Branching candidates may require piecewise affine models. Their adapter must
prove coverage of the case domains and the denotation theorem in each case,
respecting any case priority or overlap. A test basis for one branch cannot be
used as coverage of the whole function. This is particularly relevant to
Leant's explicit finite-spine case policy: enabling a structural case form is
not itself a proof of a universal affine behavior law.

A failed completion refutes only that completion. Pruning an incomplete sketch
requires a theorem that every admissible completion violates the requested
contract, or a sound impossibility result for the entire remaining search
space. One counterexample to one filled-in program is not such a theorem.

\section{Two complete mathematical interfaces}\label{sec:interfaces}

\subsection{Exact Frey quotients with separated premises}

Set
\[
 N_2=b^p-1-a^p,\qquad N_4=-a^pb^p,
\]
for $a,b\in\Z$ and $p\in\N$. The following proposition isolates the
arithmetic used by the curve-base-change proof and slightly sharpens the
sufficient lower bound needed by its first coefficient.

\begin{proposition}[Separated exactness requirements]\label{prop:frey}
If $2\mid b$, $p\ge2$, $p$ is odd, and $a\equiv3\pmod4$, then $4\mid N_2$.
If $2\mid b$ and $p\ge4$, then $16\mid N_4$, with no residue or oddness
assumption on $a$ or $p$. In each case, division by the indicated denominator
in $\Z$ gives a quotient whose image in $\Q$ equals the corresponding rational
quotient.
\end{proposition}

\begin{proof}
Write $b=2c$. If $p\ge2$, then $4\mid b^p$. Since $a\equiv-1\pmod4$ and
$p$ is odd, $a^p\equiv-1\pmod4$. Therefore
\[
 N_2\equiv0-1-(-1)=0\pmod4.
\]
If $p\ge4$, then $16\mid b^p$, and multiplication by $-a^p$ preserves that
divisibility. For either numerator $N$ and denominator $d$, exactness gives
$N=dq$ for the integer quotient $q$. Mapping to $\Q$ yields $N=dq$ there;
as $d\ne0$ in $\Q$, division gives the desired rational identity.
\end{proof}

The lower bound for $N_2$ is a sufficient condition for this route, not a claim
of a weakest possible contract. For particular $a,b,p$, exactness can hold
under weaker assumptions. The operation's fundamental guard is simply the
needed divisibility; the displayed arithmetic hypotheses are one registered
route to it.

The proposed use site is:
\begin{lstlisting}
fix a b : Int
fix p : Nat where p >= 4 and odd p
assume a == 3 modulo 4
assume 2 divides b

let q2 := exact_quotient(b^p - 1 - a^p, 4)
let q4 := exact_quotient(-a^p * b^p, 16)

know cast(q2, Rational) = (b^p - 1 - a^p)/4
  by transport_exact_quotient
know cast(q4, Rational) = (-a^p * b^p)/16
  by transport_exact_quotient
\end{lstlisting}
The rational expressions in the conclusions are elaborated over $\Q$ with the
integer inputs explicitly embedded. Their meanings are fixed before the
arithmetic profile runs. The expanded obligation graph records the first
coefficient's residue and oddness route separately from the second coefficient's
power-divisibility route.

The joint fixture $(a,b,p)=(3,2,5)$ gives $q_2=-53$ and $q_4=-486$. Removing
one of the common sufficient premises has instructive neighboring cases:
$(3,2,4)$ fails the first coefficient's exactness; $(3,3,5)$ fails it as well;
$(3,2,3)$ fails the second coefficient's exactness; and $(1,2,5)$ fails the
first. These are executed exact-arithmetic checks in the companion, while the
general proposition above is a written proof.

A reusable helper should be exported with only its stated arithmetic premises,
or checked in a deliberately restricted local context. Importing a full
contradictory counterexample package would allow an irrelevant proof by
contradiction. An assumption-support report helps detect such a bypass, but
establishing satisfiability of the helper's premises with an explicit fixture
is an additional and useful test.

\subsection{What survives a general ring map}

For a ring homomorphism $\phi:\Z\to R$, exactness always yields
\[
 \phi(N)=\phi(d)\phi(q).
\]
Turning this balance into a division formula requires the target-side
hypotheses of the selected operation. In a field, nonzeroness of $\phi(d)$
is enough. In a general ring, a chosen unit is the relevant interface.
Nonzeroness of $d$ in the source is insufficient: its image may be zero in a
positive-characteristic target.

This example motivates a useful Clover rule: a transport method exposes both
what is preserved unconditionally and the additional premises needed by its
requested consumer. It may propose the balance equation as an alternative,
but it cannot silently deliver that weaker result when the source asked for
a quotient identity.

\subsection{A representing algebra as a dependent construction}

For the finite tensor example, the natural Clover result is a dependent package:
\begin{lstlisting}
construct W : CommutativeAlgebra A
  with finite_module W
  with free_module W
  with e : for each T : CommutativeAlgebra A,
             AlgHom(W,T) equivalent_to family i, AlgHom(H(i),T)
  with natural e in T
\end{lstlisting}
The declared input interface must retain the finite and free module properties
of every $H_i$. For a finite enumerated family, an iterated tensor product
supplies the object; without a chosen enumeration, a mathematical existence
interface can follow the source's finite-type induction. These are different
construction policies with related mathematical outputs.

The naturality field has an actual equation:
\[
 e_{T'}(u\circ f)=\bigl(i\mapsto u\circ e_T(f)_i\bigr)
\]
for every $A$-algebra map $u:T\to T'$ and $f:W\to T$. The construction is
not accepted merely because each individual $e_T$ is bijective. Clover's
obligation telescope preserves this universally quantified compatibility law.

The missing-factor negative case is especially valuable. With one factor,
representability determines $W$ up to an algebra isomorphism with that factor.
Taking $A=\Q$ and $H=\Q[X]$ shows that finite indexing alone cannot force $W$
to be a finite $A$-module: the monomials give arbitrarily large linearly
independent finite families. A compact surface that dropped the factor's
module hypothesis would be expressing a false theorem, not merely a more
convenient version of the source theorem.

Finally, uniqueness up to isomorphism does not authorize substituting a different
representing object into dependent data without transport. The selected $W$
and selected equivalence family remain explicit construction results. Their
proof fields may be implicit at use sites; their data identity may not.

\section{Certified computation inside the same discipline}\label{sec:cas}

\subsection{The CAS returns a relation to the original problem}

Clover treats a computer-algebra operation as a producer of a value and evidence
for a fixed relation. Useful result specifications include exact equality,
equality under guards, equality on a region, a finite jet, a solution witness,
a complete solution characterization, and an enclosure with an error bound.
They are not points on a single ladder of confidence.

For example, an algorithm may produce $y$ with $F(y)=0$. That proves existence
of a root, not completeness of a returned list of roots. A root-coverage theorem
without a soundness theorem for listed candidates is also insufficient for a
complete solution set. A polynomial residual modulo $X^N$ can identify a jet
only with the additional hypotheses of the appropriate lifting theorem,
including the chosen constant term and an invertibility condition where needed.

The same distinction governs familiar radicals. Over the reals,
$\sqrt{x^2}=|x|$ holds for all $x$, while $\sqrt{x^2}=x$ requires $x\ge0$.
A symbolic simplifier that returns the latter expression must either supply the
guard or preserve the absolute value. The author-facing property system makes
such domain and branch information reusable, but does not turn a CAS convention
into a theorem.

\subsection{Checking and reification have different soundness obligations}

A certificate checker can have an excellent theorem about an expression
language while still be attached to the wrong original target. Clover therefore
requires both a checker theorem and a reification bridge.

For a ring-expression certificate, the bridge fixes the coefficient ring,
number of atoms, and valuation assigning each atom its exact Lean expression.
It produces equalities between the original typed expressions and the
denotations of their reified forms. The checker theorem proves the desired
relation between those denotations. Composition of these proofs establishes
the original request.

A basic ideal-membership certificate has the form
\[
 p=\sum_{i=1}^s q_i g_i.
\]
After verifying the identity in the correct polynomial ring, evaluation at a
common zero of the $g_i$ proves $p=0$ there. A certificate of $p^k$ in the ideal
is a different claim: concluding $p=0$ can require a reducedness or field
hypothesis. A certificate over $\Q[X]$ does not automatically prove the
corresponding ideal membership over $\Z[X]$.

Arity checks precede pairing of coefficient vectors, atoms, generators, or
certificate components. A truncated \code{zip} is not a proof that omitted
components are irrelevant. Every bridge binds the complete original expression,
not just its normalized printed spelling.

\subsection{A shared reifier framework, not an untyped universal normalizer}

Several checkers can share a typed reification framework and source-to-expression
proof construction. They need not share one untyped arithmetic universe.
Natural lengths, integer polynomials, rational affine forms, and real-valued
inequalities have different operations and embeddings.

The reusable component should therefore be indexed by its algebraic signature,
carrier, and named embeddings. It provides expression formation, atom identity,
and compositional denotation lemmas. An affine normalizer, polynomial checker,
and order certificate then supply different soundness theorems over appropriate
instances of that framework.

For inequalities, coefficient signs and denominator directions matter. Mapping
a natural length into $\Q$ can make affine comparison convenient, but returning
a statement about natural lengths requires the corresponding injectivity or
order-reflection theorem. The framework records that route rather than erasing
the distinction between the source and target number systems.

\subsection{Verified algorithms and checked certificates}

Clover admits two principal routes. A verified algorithm computes its result
inside a proved semantic framework. Alternatively, an untrusted producer returns
a certificate checked by a verified checker. Existing proof-producing tactics
are another provider route whose generated proof term is checked in the usual
way.

A theorem about a checker does not verify an arbitrary native execution of that
checker. The strict route retains a term whose decisive computations are checked
by kernel reduction or reconstructs a proof. Current Lean documentation records
native evaluation through additional axiom dependencies, including per-invocation
axioms for \code{native_decide}; Clover must audit the actual pinned toolchain
rather than assume an older mechanism \cite{lean-axioms}.

There is no measured certificate-size threshold in this article at which a
native route becomes preferable. Such a threshold would depend on the checker,
certificate distribution, hardware, proof policy, and retained evidence format.
The default proof policy is chosen before such optimization.

\subsection{Effectful programs are a separate contract axis}

For stateful or effectful programs, pure input/output predicates are insufficient.
Clover should use a specified Hoare-style semantics, with explicit state,
exceptions, resources, and a distinction between partial and total correctness.
The relation to weakest-precondition transformers is useful here; Dijkstra
monads provide a relevant account of such specification structure
\cite{dijkstra}.

Lean's \code{mvcgen} is an existing verification-condition-generation component,
not a feature Clover needs to reinvent. The round-three reports' concrete
Lean 4.32 probe described that interface as experimental; a production adapter
must be pinned and tested \cite{mvcgen,syn-b}. A synthesis timeout is an
operational event, never a premise inside a Hoare theorem.

\section{Should Lean's kernel type system change?}\label{sec:kernel}

\subsection{Yes to richer source typing; distinguish the kernel question}

The user-facing proposal does extend the typing discipline: operations request
proved properties, functions have behavioral interfaces, and unfinished terms
carry dependent obligations. All three can be translated into existing Lean.
This is the recommended first implementation, not an assertion that kernel
research is uninteresting.

There are two different motivations for a kernel extension. One is to express
properties that existing types cannot express. The properties in this article
do not supply such a motivation: predicates, dependent functions, subtypes, and
structures already express them. The other is to make common coercions and
transport behavior definitionally simpler or operationally cheaper. That is a
real but measurable research question.

\subsection{A primitive proof-explicit refinement former}

A candidate primitive could have the following rules for $A:\Type_u$ and
$P:A\to\Prop$:
\[
 \frac{a:A\qquad p:P(a)}{\mathsf{pack}(a,p):\CRef(A,P)},
 \qquad
 \frac{r:\CRef(A,P)}{\mathsf{forget}(r):A},
\]
with
\[
 \mathsf{forget}(\mathsf{pack}(a,p))\ \longrightarrow\ a.
\]
A subsumption operation would still require evidence:
\[
 \mathsf{weaken}:\bigl(\forall x,P(x)\to Q(x)\bigr)
                   \to\CRef(A,P)\to\CRef(A,Q).
\]
Its straightforward translation is Lean's subtype representation. The primitive
must not allow a bare $a:A$ to be regarded as refined without constructing a
proof of $P(a)$.

Such a primitive could make a frontend implementation more uniform. It does
not, merely by existing, eliminate proof obligations or improve runtime storage.
Existing subtype representations already erase proposition-valued evidence in
compiled code, and proof irrelevance removes distinctions between proofs of
one proposition \cite{lean-subtype,lean-propositions}. A convincing performance
case must identify an overhead that remains after an optimized ordinary-Lean
encoding.

\subsection{A contract arrow is a useful surface constructor}

An explicit behavioral arrow could be written schematically as
\[
 A\xrightarrow{P;Q}B
 :=\{f:A\to B\mid\forall x,P(x)\to Q(x,f(x))\}.
\]
It has introduction from a function and law proof, projection to the function,
and proof-backed variance and composition. Dependent outputs generalize $B$ to
$B(x)$. Relational laws can be bundled alongside this arrow rather than forced
into a unary postcondition.

Again, the logical meaning is already representable. The useful source-level
extension is that applications request and propagate these interfaces without
requiring every local function variable to be rebound to a new wrapper type.
A primitive kernel arrow would need a reason beyond prettier printing.

\subsection{Definitional functoriality is a more specific experiment}

A less superficial kernel experiment concerns coercions through data
constructors. For an unknown list $xs$, equations such as
\[
 \mathsf{map}(\id,xs)=xs,
 \qquad
 \mathsf{map}(g,\mathsf{map}(f,xs))
       =\mathsf{map}(g\circ f,xs)
\]
are ordinary propositional laws, not generally obtained by reducing the unknown
list's constructors. Repeated coercions through dependent containers can make
such laws relevant to elaboration and proof size.

Laurent, Lennon-Bertrand, and Maillard study definitional functoriality and its
relationship to coercive subtyping in a dependent type theory
\cite{functoriality}. That work is a serious reference point, not a drop-in
soundness result for every feature of Lean's kernel. Clover could investigate
a restricted coercion language whose composition is normalized canonically,
first by theorem-backed elaboration and only then, if justified, by new
conversion rules.

The initial experiment should compare three implementations of exactly the
same typed coercion workload: ordinary library rewriting, a source-level
canonicalizer that emits equality proofs, and a proposed kernel conversion
extension. This separates convenience from a genuine need to change the trusted
checker.

\subsection{What the metatheory would have to establish}

A kernel extension needs precise formation, introduction, elimination, and
reduction rules, together with subject reduction, compatibility with universes
and proof irrelevance, a decidable checking procedure for its stated fragment,
and a conservativity or translation argument. A conversion extension also
needs a coherent treatment of competing reduction and coercion paths.

For example, two isomorphisms of a vector space with itself can transport a
vector differently. Identity and negation on a nontrivial real vector space
are not interchangeable merely because both are isomorphisms. A coercive
subtyping scheme that silently chooses between them changes data, not just
proofs. Coherence must therefore concern the actual chosen maps, not a slogan
that isomorphic objects are ``the same''.

Lean4Lean offers relevant work toward a verified account of Lean's checking
infrastructure, but it should not be cited as an already established proof of
an unrelated kernel extension \cite{lean4lean}. The extension must be checked
against the actual version and supported feature set.

\subsection{Why semantic entailment should not enter conversion by default}

A rule identifying refined types whenever their predicates are logically
equivalent would make type conversion depend on theorem proving unless the
predicate fragment were carefully restricted. The objection is not that a
timeout becomes false; a sound implementation would return unknown. The
problem is that the foundational checking boundary would acquire an additional,
potentially expensive and incomplete logical search dependency.

Clover can obtain most of the intended convenience by inserting explicit
proof-directed coercions in the elaborator. The kernel then checks the
resulting term rather than re-running the search. A deliberately restricted,
decidable predicate-conversion fragment remains a possible experiment, but
its grammar, decision procedure, and interaction with dependent types must be
part of the proposal, not left implicit.

The practical decision is therefore: implement the richer author-facing type
system now over ordinary Lean; retain proof-explicit native refinements and
restricted definitional coercions as testable research options; move no open-ended
property inference into kernel conversion.

\section{The executable companion: what was actually built}\label{sec:prototype}

\subsection{An intentionally small source language}

The companion implements endomorphisms of one arbitrary carrier and four
properties: injectivity, involution, left inverse, and pointwise equality.
It has function binders, assumptions, claims, nested lexical scopes, and written
composition expressions. Its actual syntax is:
\begin{lstlisting}
fix f g r k : End
assume retraction : left_inverse r f
assume outer : injective g
claim composite : injective (g . f)
assume renamed : same g k
claim transported : injective (k . f)

scope
  fix j : End
  assume local : involutive j
  claim local_composite : injective (j . f)
end

claim unsupported_strengthening : involutive k
\end{lstlisting}
The first three claims obtain checked proof plans. The last remains unresolved:
no rule or premise establishes that $k$ is an involution. It is not marked
false. A derived claim can supply evidence later; an unresolved claim cannot.

The symbol \code{.} in this restricted grammar means composition with the outer
function first, so \code{(g . f)} denotes $g\circ f$. All functions are
endomorphisms of one carrier; dependent types, heterogeneous arrows, refinements
of arbitrary Lean expressions, and the richer source notation in the article
are not implemented by this parser.

\subsection{Seven semantic rules and a separate plan checker}

The fixed registry contains involution-to-left-inverse, left-inverse-to-injective,
composition of injective maps, symmetry and transitivity of pointwise equality,
transport of injectivity through pointwise equality, and reflexive pointwise
equality. Their mathematical meanings are expressed in the ordinary-Lean file
\code{CloverCore.lean}.

Search produces proof graphs. A separate checker reconstructs every rule
application, checks the instantiated premises and conclusion, verifies that
assumptions have the expected identities and propositions, checks scope, and
compares the proof with the exact request. The checker shares the expression
representation and substitution utilities with the generator; it is not an
independently verified implementation of the semantics. It does not call Lean.

The exporter produces ordinary Lean declarations with shared intermediate
\code{have} statements rather than recursively duplicating whole proofs.
For the first example, the mathematical content of the generated body is:
\begin{lstlisting}
have step1 : Injective f := inj_from_left_inverse f r retraction
have step2 : Injective (compose g f) := inj_comp f g step1 outer
exact step2
\end{lstlisting}
The actual generated identifiers include their unique local identities. The
exported theorem retains all original ambient assumptions; its separate support
report identifies the assumptions actually used. It does not silently publish
a strengthened theorem by deleting unused premises.

\subsection{Executed regression results}

The 35 unit-test methods pass. They cover positive composition and transport,
seedless cycles, missing premises, resource exhaustion, changed snapshot and
target identities, forged substitutions, forbidden rules, a required root
rule, sibling-scope isolation, shadowed names, pending claims, malformed
syntax, and ordinary Lean export. One method compares opposite closure
schedules across 32 seed subsets. Another exhausts 136 valuations of the seven
rules over all functions on a two-element carrier.

Those finite semantic checks are not universal proofs of the rules. The
ordinary mathematical proofs of the rules are simple, and their Lean reference
contains no \code{sorry} or added axioms in its source, but it was not compiled
in this environment. Neither the source inspection nor the Python checks are
reported as kernel acceptance.

A separate exact-arithmetic program checks the affine test-basis construction,
its concrete realizers, 625 small affine forms, finite grids for accepted forms,
empty-element-type and guarded-length examples, and the Frey fixtures.
Its report records 1,576 individual assertion checks. This count is not added
to the number of unit-test methods as though they were the same unit of evidence.

\subsection{A small but informative profile experiment}

\begin{center}\small
\begin{tabularx}{\textwidth}{@{}Yrr@{}}
\toprule
Metric, same first composition request & Seven rules & Two-rule profile\\\midrule
Arena terms & 5 & 5\\
Candidate instantiations / retained ground rules & 211 & 16\\
Relevant ground atoms & 60 & 18\\
Rule tests in repeated-scan closure & 633 & 48\\
New closure facts & 7 & 2\\
Nodes in retained proof, including assumptions & 4 & 4\\
\bottomrule
\end{tabularx}
\end{center}

The narrow profile permits only left-inverse-to-injective and injective
composition. The reference implementation uses repeated scans, not the optimized
premise-counter closure described earlier. Its counters therefore expose both
grounding cost and avoidable closure work. The experiment supports a specific
engineering decision---make profiles small and inspectable---not a claim that
Clover is faster than Lean, \code{fun_prop}, or direct theorem lookup.

\subsection{What these tests do not cover}

The companion does not implement a real Lean local-context index, universe
metavariables, instance selection, dependent transport, analytic locality,
Frey property-implicit calls, full behavioral synthesis, a CAS certificate
bridge, or incremental editor migration. It does not compile the generated
Lean files. It does not run the full round-three negative suite, reproduce the
syntheses' reported kernel experiments, or measure user comprehension.

Its advance is narrower: a real parser and executable finite grounding policy
produce scope-sensitive proof plans and ordinary Lean source, and the test suite
mutates the same artifacts at the correspondence boundaries. The next useful
step is to port this pipeline to Lean and replace the symbolic plan-checking
boundary with actual kernel-checked terms, not to enlarge the Python model
until it becomes an unofficial second Lean.

\section{A Lean implementation plan with explicit exit tests}\label{sec:implementation}

\subsection{Modules and responsibilities}

The first Lean package should have a small separation of concerns.
\code{Clover.Syntax} defines the structured source nodes and source spans.
\code{Clover.Seal} classifies holes and retains the original typed targets.
\code{Clover.Registry} validates theorem-role metadata.
\code{Clover.Evidence} indexes local proofs and quantified schemes.
\code{Clover.Obligations} stores typed pending requests and dependency edges.
\code{Clover.Resolve} implements lookup, bounded grounding, and profile selection.
\code{Clover.Export} reconstructs accepted terms and source-to-proof mappings.
A renderer consumes these records rather than introducing another semantics.

Provider adapters are separate modules. They may call existing tactics or Leant,
but they do not mutate the accepted evidence index directly. They return candidate
results to a commit boundary that checks the exact target, dependencies, and
policy before admitting any evidence.

\subsection{Gate one: replace the symbolic boundary with Lean}

Port the small endomorphism example first. The gate is not merely that the seven
library theorems compile; it is that source claims elaborate into checked Lean
terms and the same negative mutations fail at the intended boundary. A missing
premise must leave an obligation. A scope leak must be rejected. A forged result
for another target must not be accepted because it is a true theorem.

The Lean implementation should log the actual theorem instances, number and
size of retained expression nodes, grounding work, and kernel-checking outcome.
The Python implementation can serve as a reference for the small finite grammar,
but disagreements must be resolved semantically rather than by treating Python
as an oracle for Lean.

\subsection{Gate two: one property-implicit mathematical call}

Implement \code{exact_quotient} and its transport interface over the existing
Lean/Mathlib arithmetic lemmas. Reuse or extend the syntheses' already reported
Frey artifacts rather than rebuilding a broad arithmetic library. Exercise the
satisfiable fixture and the separated coefficient requirements from
Proposition~\ref{prop:frey}.

The exit test compares a property-implicit call with a normal Lean call supplied
with the same helper theorems and automation. If the new interface adds
unpredictable search while removing only a few argument names, the result may
belong in a library or automatic-parameter helper rather than a new language.

\subsection{Gate three: quantified locality}

Implement one regional-equality scheme and the consumer turning it into a
neighborhood equality at a member of an open set. Use the autonomous-derivative
proof without strengthening its attained-range differentiability assumptions.
The negative tests remove openness, specialize the induction hypothesis too
early, or replace neighborhood equality by equality at one point.

This gate tests something the endomorphism model cannot: binder scope,
quantified theorem reuse, theorem-role matching, and proof transport through
actual analytic interfaces. Existing \code{fun_prop} rules and ordinary Lean
wrappers are controls, not competitors to be excluded from the comparison
\cite{fun-prop}.

\subsection{Gate four: a behavioral denotation bridge}

Port a small admitted list grammar and its length soundness proof into the
Leant adapter. The required end-to-end result is a proof about the exact
candidate's original contract. Then attach a realizable affine test basis,
including the correlated-pair and empty-element cases.

The gate fails if the result is only a successful model check with no source
bridge, if a passive provider assertion is treated as a Lean law, or if a
counterexample lacks a source input satisfying the declared guard. Piecewise
models need explicit case coverage. A broad solver connection is not a
substitute for this narrow correspondence proof.

\subsection{Gate five: replay, repair, and publication}

Add snapshot-local invalidation and typed transport for one selected class of
rewrites. Test changes of binder identity, structure dictionary, measure,
universe instantiation, dependent index, and method policy. Initially rebuild
rather than migrate whenever correspondence is uncertain.

Publication requires two separate checks. Every asserted document node must
have an accepted proof or be visibly marked unfinished; and every exported
mathematical result must meet its pinned axiom policy. A failed unused assertion
blocks a fully verified document even when another route proves the final
theorem. An ordinary partial document may still be saved and inspected; it must
not be mislabeled as completely checked.

\section{Evaluation that can disconfirm the language proposal}\label{sec:evaluation}

\subsection{Separate infrastructure from notation}

Use five principal experimental arms. $L_0$ is stock Lean and the project's
existing tools. $L_1$ adds the new mathematical packages and theorem registrations
but uses ordinary Lean syntax. $L_2$ adds the same search, checking, and certificate
services available to Clover. $L_3$ adds property-aware use sites and the shared
obligation engine. $L_4$ adds the proposed mathematical surface. A renderer-only
control isolates the benefit of a more readable view without a new input
language.

Packages, wrapper lemmas, and relevant existing automation must be available
consistently across the appropriate arms. A one-line Clover invocation is not
compared with an unassisted Lean proof that has been denied the theorem it calls.
The most informative comparisons are $L_3$ versus $L_2$, and $L_4$ versus
$L_3$ plus the renderer control.

\subsection{Measure semantic work, not just characters}

Primary authoring measures should include time spent actively constructing and
repairing proofs, interventions needed to supply missing meaning, and incorrect
interpretations of compact statements. Machine measures should separate
elaboration, grounding, external search, proof construction, kernel checking,
proof-term size, and retained index memory. Report distributions and outliers,
not only a mean.

A reading experiment should ask participants to reconstruct the domain,
quantifier order, hypotheses, conclusion strength, and claimed proof method
before seeing the expanded view. The almost-everywhere quantifier example,
exact quotient cast, and finite tensor construction are suitable development
items because their omitted distinctions can be stated precisely. They are
not held-out items after they have influenced this design.

All examples reviewed in this article belong to the development set. Held-out
tasks should be selected after the interfaces are fixed, include already-short
Lean proofs, and include cases where the intended inference fragment fails.
Participants and task order should be counterbalanced so that familiarity with
one source does not masquerade as a language benefit.

\subsection{Negative tests need a positive neighbor and a failure kind}

Every semantic mutation should record a satisfiable positive fixture where
applicable, the changed premise or target, and the expected failure category.
Some representative cases are:
\begin{center}\small
\begin{longtable}{@{}P{.28\textwidth}P{.42\textwidth}P{.20\textwidth}@{}}
\toprule
Positive situation & Mutation & Required outcome\\\midrule\endfirsthead
\toprule Positive situation & Mutation & Required outcome\\\midrule\endhead
\bottomrule\endfoot
A left inverse of the same function & Rename or replace that function at the consumer & Correspondence failure or unresolved property.\\
Regional equality near a point & Keep only equality at the point & Locality obligation remains.\\
An exact integer quotient & Remove divisibility, keeping the same rational-looking notation & Exactness obligation remains.\\
A target-side unit & Keep only source-side nonzeroness & Target inverse obligation remains.\\
A finite family of finite free factors & Remove finite generation of one factor & Construction contract is not established.\\
Countably many a.e. equalities & Replace the family by an arbitrary uncountable one & Uniform-null-set obligation remains.\\
A concrete length counterexample & Use an unrealizable or unguarded observation & No source refutation.\\
A law about one candidate & Change its term graph or displayed edit & Recheck; old receipt is insufficient.\\
A completed proof document & Insert an unused unsupported assertion & Document is not fully checked.\\
A method-selected derivation & Supply a different proof route without the requested method & Method-policy failure, not mathematical falsity.\\
\end{longtable}
\end{center}

Some failures are logical counterexamples; others are deliberately outside an
automatic fragment or violate correspondence policy. The test runner must not
collapse them into a single boolean meaning ``false.'' The companion executes
only the subset described in Section~\ref{sec:prototype}; this table is the
larger Lean implementation's acceptance specification.

\subsection{Stopping rules}

If the library and service arms achieve the same authoring improvement as the
property layer, retain the library, providers, and renderer. If the property
layer helps but the new syntax harms comprehension or repair, retain the
property-aware Lean extension without a separate language. If a small profile
works well but broad automatic inference is unpredictable, keep the small
profile and explicit method selection.

These are not retreats from the project. They are the distinctions needed to
identify which part of the design actually solves the observed problem.
No human productivity result or Lean elaboration speedup is claimed in this
article.

\section{Answers to the round-three implementation questions}\label{sec:answers}

The syntheses' questions overlap. The following table gives a concrete disposition
for the most important shared demands and the ten explicitly numbered questions
in \emph{Nine Refinements}, without pretending that a design answer is a
completed experiment.

\begin{center}\small
\begin{longtable}{@{}P{.28\textwidth}P{.63\textwidth}@{}}
\toprule
Question or demand & Clover's answer and current status\\\midrule\endfirsthead
\toprule Question or demand & Clover's answer and current status\\\midrule\endhead
\bottomrule\endfoot
T1: cost on a real Lean file & Still unmeasured. The companion reports model operation counts; the implementation gates require separate real Lean timings, proof sizes, and index memory.\\
T2: which rules fire? & Retain exact proof-DAG rule instances and all grounding counters. The first model proof uses two rules; the full seven-rule profile generates substantial unused search. Real-file comparison with direct lookup remains open.\\
T3: rewritten anchors & Preserve the original anchor and use checked equality or proposition-equivalence bridges. Weaker relations require consumer theorems. General Lean rewrite migration is not implemented here.\\
T4: a common reifier & Share a typed signature-indexed framework and denotation construction, not one untyped arithmetic grammar. Original-target bridges remain checker-independent obligations.\\
T5: abstraction completeness & Theorem~\ref{thm:basis} supplies a complete affine criterion for an observation image with coverage and realizers. It does not establish a bridge for every current Leant Length feature.\\
T6: almost-everywhere rows & Store the measure and full binder order. Register integral consumers and countable-family collection explicitly; point evaluation is not a default consumer.\\
T7: lexical structure contexts & Fix actual dictionaries and create local instance bridges at use sites. No generated-preamble shortening experiment was performed.\\
T8: bounded rule generation & Section~\ref{sec:grounding} specifies source-derived arenas and goal-triggered grounding; the companion implements this for a fixed endomorphism grammar.\\
T9: reader recovery & Use blind reconstruction of domain, quantifiers, hypotheses, conclusion strength, and method. No reading study has been conducted.\\
T10: independent implementations & The symbolic checker and search engine are separate paths but share representations and are not independent Lean implementations. A Lean port compared against the model is the next useful comparison.\\
Witness choice and method fidelity & Classify holes by semantic role; permit new data only at construction nodes. Check an explicit derivation skeleton or declared rule policy rather than a decorative method name.\\
Trusted execution and migration & Retain proof/certificate routes under a pinned axiom policy. Recheck after environment changes; no measured native-evaluation crossover is claimed.\\
First interface and exit condition & Start with the narrow checked pipeline, then exact quotients and regional differentiation. Keep only the library or renderer if controlled comparisons do not support the language layer.\\
\end{longtable}
\end{center}

\section{Conclusion}

Clover's central claim is not that mathematical rigor needs fewer obligations.
It is that authors should be able to state the mathematical scope and
transformation that generate those obligations, while a disciplined elaborator
assembles the routine evidence.

The round-three refinement calculus is a suitable foundation. The next step
is to make its operational boundary concrete: retain quantified evidence;
represent unfinished results as genuinely conditional terms; restrict routine
inference to explicit finite fragments; preserve object identity through
rewriting and replay; and equip behavioral observations with the coverage and
realization theorems their use requires.

The executable slice demonstrates part of that boundary and exposes a cost
that a purely architectural proposal could hide. The affine test-basis theorem
shows how richer behavioral typing can improve both positive verification and
negative witness validity without extending kernel conversion. The source
studies show where the same discipline matters in analysis, arithmetic, and
algebraic construction.

The recommended Clover implementation is therefore a small proof-producing
extension of Lean, evaluated against equally equipped ordinary Lean, with a
mathematical surface only where it earns its place. Native refinements or
restricted definitional coercions remain precise research options, not
prerequisites for beginning that implementation.

\appendix
\section{Reproducing the companion}\label{app:reproduce}

The package uses Python's standard library only. From the \code{companion}
directory, run:
\begin{lstlisting}
python -m unittest -v test_clover_slice
python clover_slice.py demo.clover --json demo-result.json \
  --lean GeneratedExamples.lean
python affine_basis.py
\end{lstlisting}
The first command runs the 35 unit-test methods. The second re-creates the
three accepted proof plans, the unresolved claim, and the ordinary-Lean export.
The third re-creates the exact-arithmetic report. The saved profile comparison
is reproducible by the command in the README.

To check the Lean reference and generated examples in a suitable installed
Lean environment, make the companion directory available on Lean's import
path, compile \code{CloverCore.lean} to \code{CloverCore.olean}, and then check
\code{GeneratedExamples.lean}. The README gives platform-specific shell
commands. No successful execution of those commands is claimed by this package.
The intended first target for integration is the campaign's Lean 4.32.0
baseline, followed by an explicitly pinned upgrade test.

The article is self-contained \LaTeX{} and builds with:
\begin{lstlisting}
latexmk -pdf -interaction=nonstopmode -halt-on-error Clover.tex
\end{lstlisting}
No external bibliography processor, downloaded font, or generated table file
is required.

\section{A minimal accepted-result record}\label{app:record}

The full implementation's proof-bearing result record should preserve at least
the following information. This is a schema specification, not a claim that
the symbolic companion implements every field.
\begin{lstlisting}
AcceptedStep {
  sourceNode      : SourceSpanAndStableNodeIdentity
  originalTarget  : TypedLeanExpr
  context         : DependentLocalContextSnapshot
  structureArgs   : ExplicitStructureArguments
  construction    : OptionalAuthorizedWitnessRecord
  derivation      : TypedMethodOrRuleDAG
  evidence        : LeanProofOrCheckedCertificateBridge
  conclusion      : TypedLeanExpr
  support         : TransitiveAssumptionAndDeclarationDependencies
  policy          : RuleAndAxiomPolicy
  environment     : PinnedToolchainAndImportManifest
  replay          : CheckedCorrespondenceToDisplayedSourceEdit
}
\end{lstlisting}
Hashes can index and bind serialized records, but cannot replace typed
correspondence checks. A record with an unfinished obligation is a pending
record with a different status, not an \code{AcceptedStep} with an optimistic
flag. An external producer may report a candidate record; only the checked
commit boundary constructs the accepted form.

\section{Evidence inventory}\label{app:evidence}

\begin{center}\small
\begin{tabularx}{\textwidth}{@{}P{.29\textwidth}Y@{}}
\toprule
Artifact & Evidence provided\\\midrule
\code{Clover.tex}, \code{Clover.pdf} & Design, written proofs, source analysis, explicit limits, and evaluation specification.\\
\code{SOURCE_MANIFEST.json} & Exact upstream repository references and scope of direct inspection.\\
\code{clover_slice.py} & Executable restricted parser, finite grounding, proof-plan checking, and Lean export.\\
\code{test_clover_slice.py} & 35 passing unit-test methods; some contain multiple finite cases.\\
\code{CloverCore.lean} & Ordinary-Lean definitions and proofs for the seven rules; not compiled in this environment.\\
\code{GeneratedExamples.lean} & Three exported theorem bodies and a comment marking the unresolved claim; not compiled here.\\
\code{demo-result.json} & Exact model verdicts, assumption support, and counters; all explicitly say kernel checking was not performed.\\
\code{profile-comparison.json} & Reproducible operation-count comparison on the same model request.\\
\code{affine_basis.py}, \code{affine-result.json} & Exact arithmetic, realizers, and 1,576 finite checks; not a verified source-denotation bridge.\\
\code{EXECUTION_STATUS.json}, \code{test-output.txt} & Execution scope and saved test output.\\
\bottomrule
\end{tabularx}
\end{center}

\clearpage
\begin{thebibliography}{99}
\small\raggedright
\bibitem{syn-a}
Brainstorm campaign, \emph{From Properties to Proof Obligations}, revised
round-three synthesis, commit \code{bf3333905995060870f8585e297ffc16f3fe7e86}.
Main \LaTeX{} report inspected on 6 September 2026.
\url{https://github.com/VladimirReshetnikov/Brainstorm/blob/bf3333905995060870f8585e297ffc16f3fe7e86/docs/round-3/synthesis/unified-report.tex}.

\bibitem{syn-b}
Brainstorm campaign, \emph{Nine Refinements}, revised round-three unified report,
same Brainstorm commit. Main \LaTeX{} report inspected on 6 September 2026.
\url{https://github.com/VladimirReshetnikov/Brainstorm/blob/bf3333905995060870f8585e297ffc16f3fe7e86/docs/round-3/unified_report/unified_report.tex}.

\bibitem{autonomous}
Vladimir Reshetnikov et al., ProveIt,
\emph{AutonomousIteratedDeriv.lean}, commit
\code{24ce8bd743eaab64a91ce90725ea00f498d319d2}.
\url{https://github.com/VladimirReshetnikov/ProveIt/blob/24ce8bd743eaab64a91ce90725ea00f498d319d2/Analysis/FabiusFunction/Lean/FabiusFunction/AutonomousIteratedDeriv.lean}.

\bibitem{frey}
Anthropic FLT repository, \emph{S\_FreyPackage\_freyCurveInt\_map.lean},
commit \code{aa2d8b34692b16c70f699536de0d8e75b9a3e9ef}.
\url{https://github.com/anthropics/fermats-last-theorem/blob/aa2d8b34692b16c70f699536de0d8e75b9a3e9ef/P2M/Sol/S_FreyPackage_freyCurveInt_map.lean}.

\bibitem{tensor}
Anthropic FLT repository, \emph{S\_Algebra\_exists\_algHom\_equiv\_pi.lean},
same FLT commit.
\url{https://github.com/anthropics/fermats-last-theorem/blob/aa2d8b34692b16c70f699536de0d8e75b9a3e9ef/P2M/Sol/S_Algebra_exists_algHom_equiv_pi.lean}.

\bibitem{leant-verification}
Vladimir Reshetnikov et al., Leant, \emph{src/Leant/Synth/Verification.hs},
commit \code{3a40904be8a410d832d9ab6900b3d3e7b425eccb}.
\url{https://github.com/VladimirReshetnikov/Leant/blob/3a40904be8a410d832d9ab6900b3d3e7b425eccb/src/Leant/Synth/Verification.hs}.

\bibitem{leant-contract}
Leant, \emph{src/Leant/Synth/Length/Contract.hs}, same directly inspected commit.
\url{https://github.com/VladimirReshetnikov/Leant/blob/3a40904be8a410d832d9ab6900b3d3e7b425eccb/src/Leant/Synth/Length/Contract.hs}.

\bibitem{leant-adapter}
Leant, \emph{src/Leant/Synth/Length/Adapter.hs}, same directly inspected commit.
\url{https://github.com/VladimirReshetnikov/Leant/blob/3a40904be8a410d832d9ab6900b3d3e7b425eccb/src/Leant/Synth/Length/Adapter.hs}.

\bibitem{lean-subtype}
Lean Language Reference, \emph{Subtypes}, online documentation accessed
6 September 2026.
\url{https://lean-lang.org/doc/reference/latest/Basic-Types/Subtypes/}.

\bibitem{lean-propositions}
Lean Language Reference, \emph{Propositions}, accessed 6 September 2026.
\url{https://lean-lang.org/doc/reference/latest/The-Type-System/Propositions/}.

\bibitem{lean-quantifiers}
Lean Language Reference, \emph{Quantifiers}, accessed 6 September 2026.
\url{https://lean-lang.org/doc/reference/latest/Basic-Propositions/Quantifiers/}.

\bibitem{lean-application}
Lean Language Reference, \emph{Function Application}, accessed 6 September 2026.
\url{https://lean-lang.org/doc/reference/latest/Terms/Function-Application/}.

\bibitem{lean-synthetic}
Lean API documentation, \emph{Lean.Elab.SyntheticMVars}, accessed
6 September 2026.
\url{https://lean-lang.org/doc/api/Lean/Elab/SyntheticMVars.html}.

\bibitem{lean-elaborators}
Lean Language Reference, \emph{Elaborators}, accessed 6 September 2026.
\url{https://lean-lang.org/doc/reference/latest/Notations-and-Macros/Elaborators/}.

\bibitem{lean-axioms}
Lean Language Reference, \emph{Axioms}, especially the parametricity example,
native-evaluation discussion, and dependency audit, accessed 6 September 2026.
The \code{latest} documentation observed here contains links to the
4.34.0-rc2 reference; it is not used as a substitute for the campaign's
4.32.0 test pin.
\url{https://lean-lang.org/doc/reference/latest/Axioms/}.

\bibitem{fun-prop}
Mathlib documentation, \emph{Mathlib.Tactic.FunProp}, accessed 6 September 2026.
\url{https://leanprover-community.github.io/mathlib4_docs/Mathlib/Tactic/FunProp.html}.

\bibitem{mvcgen}
Lean Language Reference, \emph{The mvcgen tactic: Overview}, accessed
6 September 2026.
\url{https://lean-lang.org/doc/reference/latest/The--mvcgen--tactic/Overview/}.

\bibitem{liquid}
Patrick M. Rondon, Ming Kawaguchi, and Ranjit Jhala,
\emph{Liquid Types}, PLDI 2008.
\url{https://doi.org/10.1145/1375581.1375602}.

\bibitem{refinement-tutorial}
Ranjit Jhala and Niki Vazou, \emph{Refinement Types: A Tutorial},
arXiv:2010.07763, 2020.
\url{https://arxiv.org/abs/2010.07763}.

\bibitem{dijkstra}
Danel Ahman et al., \emph{Dijkstra Monads for Free}, POPL 2017;
arXiv:1608.06499.
\url{https://arxiv.org/abs/1608.06499}.

\bibitem{functoriality}
Th\'eo Laurent, Meven Lennon-Bertrand, and Kenji Maillard,
\emph{Definitional Functoriality for Dependent (Sub)Types}, extended version,
arXiv:2310.14929; ESOP 2024.
\url{https://arxiv.org/abs/2310.14929}.

\bibitem{lean4lean}
Mario Carneiro, \emph{Lean4Lean: Towards a Verified Typechecker for Lean,
in Lean}, arXiv:2403.14064, 2024.
\url{https://arxiv.org/abs/2403.14064}.
\end{thebibliography}
\end{document}
